\newcommand*{\DocType}{scrreprt}
%\renewcommand*\DocType{article} % Uncomment for screen optimization

\newcommand*\ClassList{scrreprt,article}
\documentclass[\DocType, paper=a4,fontsize=11pt,titlepage]{generalclass}


\usepackage[utf8]{inputenc}
\usepackage[automark]{scrpage2}	        % Seiten-Stil für scrartcl
% Mathematische Zeichensätze und Umgebungen
\usepackage{amsfonts, amssymb}	        % Definition einer Liste mathematischer Fontbefehle und Symbole
\usepackage[intlimits,sumlimits]{amsmath} % Integral-/Summationsgrenzen über/unter Zeichen
\usepackage{tabto}
\usepackage{mathabx}
% mathematische Verbesserungen
\usepackage{amsthm}	                    % spezielle theorem Stile
\usepackage{aliascnt} 
\usepackage{array}		                % erweiterte Tabellen
% Schriftzeichen, Format
\usepackage{latexsym}		            % Latex-Symbole
\usepackage[english, german, ngerman]{babel} % Mehrsprachenumgebung
% Layout
\usepackage{geometry}                   % Seitenränder
\usepackage{xcolor}                     % Farben
\usepackage{bbm}
% Tabellen und Listen
\usepackage{float}		                % Gleitobjekte 
\usepackage[flushright]{paralist}       % Bessere Behandlung der Auflistungen
\usepackage{datetime}
% Bilder
\usepackage[final]{graphicx}            % Graphiken einbinden
\usepackage{caption}                    % Beschriftungen
\usepackage{subcaption}                 % Beschriftungen für Unterteilung
\usepackage{tikz}
% Interaktive Referenzen, und PDF-Keys
\usepackage{xr-hyper}  
\usepackage[pagebackref,pdftex, plainpages=false]{hyperref} % Rückreferenz im Literaturverzeichnis, Treiber für ps zu pdf ; für direkt nach pdf: pdftex
\usepackage{enumitem}

\title{Nash Gleichgewichte in Dynamischen Flüssen}
\date{\todo{Get it done!}}
\author{Michael Markl}

%-------------------------------------------------------------------------------
% Kopf-Zeilen
%-------------------------------------------------------------------------------

\pagestyle{scrheadings}		     % Kopfzeilen nach scr-Standard		
\ifx\chapter\undefined 		     % falls Kapitel nicht definiert sind
  \automark[subsection]{section} % Kopf- und Fusszeilen setzen
\else				             % Kapitel sind definiert
  \automark[section]{chapter}	 % Kopf- und Fusszeilen setzen
\fi

%-------------------------------------------------------------------------------
%   Maske für Überschrift 
%-------------------------------------------------------------------------------
% Belegung der notwendigen Kommandos für die Titelseite
\newcommand{\autor}{Markl, Michael} 		% bearbeitender Student
\newcommand{\veranstaltung}{Seminar zur Optimierung und Spieltheorie} 	% Titel des ganzen Seminars
\newcommand{\uni}{Institut für Mathematik der Universität Augsburg} % Universit\"at
\newcommand{\matrikelnummer}{1474802}
\newcommand{\lehrstuhl}{Diskrete Mathematik, Optimierung und Operations Research} % Lehrstuhl
\newcommand{\semester}{Sommersemester 2019}	% Winter-/Sommersemester mit Jahr
\newcommand{\datum}{27.06.2019} 			% Datumsangabe
\newcommand{\thema}{Nash Gleichgewichte in Dynamischen Flüssen}  		% Titel der Seminararbeit

\newcommand{\ownline}{\vspace{.7em}\hrule\vspace{.7em}} % horizontale Linie mit Abstand

\newcommand{\seminarkopf}{
	% Befehl zum Erzeugen der Titelseite 
 \textsc{\autor}  \hfill{\datum} \\ 
\textbf{\veranstaltung} \\ 
\uni \\ 
\lehrstuhl \\
\semester \hfill{Matrikelnummer: \matrikelnummer}
\ownline 

\begin{center}
{\LARGE \textbf{\thema}}
\end{center}

\ownline
}			% Befehle und Pakete für Titelseite


\DeclareMathOperator{\e}{ex}
\DeclareMathOperator{\ma}{mate}
\DeclareMathOperator{\Ex}{Ex}

%-------------------------------------------------------------------------------
%   Befehle für Nummerierung der Ergebnisse
%   fortlaufend innerhalb eines Abschnittes
%-------------------------------------------------------------------------------
\theoremstyle{plain}            % normaler Stil
\newtheorem{theorem}{Theorem}
\numberwithin{theorem}{section}
% Lemma
\newaliascnt{lemma}{theorem}
\newtheorem{lemma}[lemma]{Lemma}
\aliascntresetthe{lemma}
% Satz
\newaliascnt{satz}{theorem}
\newtheorem{satz}[satz]{Satz}
\aliascntresetthe{satz}
% Korollar
\newaliascnt{corollary}{theorem}
\newtheorem{corollary}[corollary]{Korollar}
\aliascntresetthe{corollary}
% Proposition
\newaliascnt{proposition}{theorem}
\newtheorem{proposition}[proposition]{Proposition}
\aliascntresetthe{proposition}
%-------------------------------------------------------------------------------
\theoremstyle{definition}	% Definitionsstil
% Definition
\newaliascnt{definition}{theorem}
\newtheorem{definition}[definition]{Definition}
\aliascntresetthe{definition}
% Beispiel
\newaliascnt{example}{theorem}
\newtheorem{example}[example]{Beispiel}
\aliascntresetthe{example}
% Problem
\newaliascnt{problem}{theorem}
\newtheorem{problem}[problem]{Problem}
\aliascntresetthe{problem}
% Algorithmus
\newaliascnt{algorithmus}{theorem}
\newtheorem{algorithmus}[algorithmus]{Algorithmus}
\aliascntresetthe{algorithmus}
%-------------------------------------------------------------------------------
\theoremstyle{remark}		% Bemerkungsstil
% Bemerkung
\newaliascnt{remark}{theorem}
\newtheorem{remark}[remark]{Bemerkung}
\aliascntresetthe{remark}
% Vermutung
\newaliascnt{conjecture}{theorem}
\newtheorem{conjecture}[conjecture]{Vermutung}
\aliascntresetthe{conjecture}
% Notation
\newaliascnt{notation}{theorem}
\newtheorem{notation}[notation]{Notation}
\aliascntresetthe{notation}

%-------------------------------------------------------------------------------
% automatische Referenzen mit interaktiven Text
%-------------------------------------------------------------------------------

% Texte
\renewcommand{\theoremautorefname}{Theorem}
\newcommand{\lemmaautorefname}{Lemma}
\newcommand{\satzautorefname}{Satz}
\newcommand{\korollarautorefname}{Korollar}
\newcommand{\propositionautorefname}{Proposition}

\newcommand{\definitionautorefname}{Definition}
\newcommand{\beispielautorefname}{Beispiel}
\newcommand{\problemautorefname}{Problem}
\newcommand{\algorithmusautorefname}{Algorithmus}

\newcommand{\bemerkungautorefname}{Bemerkung}
\newcommand{\vermutungautorefname}{Vermutung}
\newcommand{\notationautorefname}{Notation}

%-------------------------------------------------------------------------------
% Nummerierung der Gleichungen innerhalb der obersten Ebene
%-------------------------------------------------------------------------------
\ifx\chapter\undefined 			% Kapitel sind definiert
  \numberwithin{equation}{section}	% Gleichungsnummern in Section
\else					% Kapitel sind nicht definiert
  \numberwithin{equation}{chapter}	% Gleichungsnummern in Kapiteln
\fi





\makeatletter
\newcommand{\overleftsmallarrow}{\mathpalette{\overarrowsmall@\leftarrowfill@}}
\newcommand{\overrightsmallarrow}{\mathpalette{\overarrowsmall@\rightarrowfill@}}
\newcommand{\overleftrightsmallarrow}{\mathpalette{\overarrowsmall@\leftrightarrowfill@}}
\newcommand{\overarrowsmall@}[3]{%
	\vbox{%
		\ialign{%
			##\crcr
			#1{\smaller@style{#2}}\crcr
			\noalign{\nointerlineskip\vskip0.4pt}%
			$\m@th\hfil#2#3\hfil$\crcr
		}%
	}%
}
\def\smaller@style#1{%
	\ifx#1\displaystyle\scriptstyle\else
	\ifx#1\textstyle\scriptstyle\else
	\scriptscriptstyle
	\fi
	\fi
}
\makeatother
			% Mathematische Befehle und Pakete

% Literatur-Bibliothek
\bibliographystyle{alphadin}               % deutscher Bibliotheksstil

% Erweiterte Einstellungen zu hyperref

\hypersetup{
        breaklinks=true,        % zu lange Links unterbrechen
        colorlinks=true,        % Färben von Referenzen
        citecolor=black,        % Farbe der Zitate
        linkcolor=black,        % Farbe der Links
        extension=pdf,          % Externe Dokumente können eingebunden werden.
        ngerman,		
	pdfview=FitH,
	pdfstartview=FitH,		
	bookmarksnumbered=true,     % Anzeige der Abschnittsnummern	% pdf-Titel
	pdfauthor={\autor}          % pdf-Autor
}

% Namen für Referenzen 

\newcommand{\ownautorefnames}{
  \renewcommand{\sectionautorefname}{Kapitel}
  \renewcommand{\subsectionautorefname}{Unterkapitel}
  \renewcommand{\subsubsectionautorefname}{\subsectionautorefname}
  \renewcommand{\appendixautorefname}{Anhang}
  \renewcommand{\figureautorefname}{Abbildung}
}

% Rückreferenzentext zur Literatur
\def\bibandname{und}%
\renewcommand*{\backref}[1]{}
\renewcommand*{\backrefalt}[4]{%
\ifcase #1 %
 (Nicht zitiert, also Ergänzungsliteratur.)%
\or
 (Zitiert auf Seite #2.)%
\else
 (Zitiert auf den Seiten #2.)%
\fi
}
\renewcommand{\backreftwosep}{ und~} % seperate 2 pages
\renewcommand{\backreflastsep}{ und~} % seperate last of longer 

			% Befehle und Pakete für Referenzen


\IfClass{article}{ % Optimize for screen
	\geometry{papersize={160mm,161.5mm},margin=5mm} 
}

\renewcommand{\descriptionlabel}[1]{\hspace{\labelsep}\textit{#1}}

\newcommand{\todo}[1]{{\color{red}#1}}

\numberwithin{figure}{section}	% Abbildungsnummern in Section

\newcommand{\R}{\mathbb{R}}
\newcommand{\Q}{\mathbb{Q}}
\newcommand{\Z}{\mathbb{Z}}
\newcommand{\N}{\mathbb{N}}
\newcommand*\diff{\mathop{}\!\mathrm{d}}

\renewcommand{\]}{\end{equation*}}

%%%%%%%%%%%%%%%%%%%%%%%%%%%%%%%%%%%%%%%%%%%%%%%%%%%%%%%%%%%%%%%%%%%%%%%%%%%%%%%%
% Start des Dokuments
\begin{document}

\ownautorefnames		% Änderung einiger automatischen Texte von hyperref (wie in referenz.tex definiert)

\thispagestyle{empty}
\newcommand{\Universitat}{Universität Augsburg}
\begin{titlepage}
  	\begin{center}
		\Large \textbf{\Universitat} \\
		\uni \\
		\lehrstuhl \\
	\end{center}
	\vspace*{1em}
	\begin{center}
		\includegraphics[width=5cm,height=5cm]{UNI}
	\end{center}
	
	\vspace*{1em}
	
	\begin{center}
		\Huge\textsc{Bachelorarbeit}
	\end{center}
	
	\vspace*{1cm}
	\ownline
	\begin{center}
		\begingroup\doublespacing
		{\LARGE \textbf{\thema}}\par
		\endgroup
	\end{center}
	\ownline
    \vfill
    {\setlength{\tabcolsep}{12pt}
   	\centering
	\begin{tabular}{ll}
		Betreuer: & Prof. Dr. Tobias Harks\\
		Zweitgutachter: & Prof. Dr. Mirjam Dür\\[1em]
		Vorgelegt von: & Michael Markl\\
		Matrikelnummer: & 1474802\\
		Abgabetermin: & \todo{Get it done!}
	\end{tabular}\par  
	}
\end{titlepage}


\thispagestyle{empty}
\tableofcontents        % Inhaltsverzeichnis
%\listoffigures         % Abbildungsverzeichnis (eventuell einfügen)
%\listoftables          % Tabellenverzeichnis (eventuell einfügen)
\setcounter{page}{0}    % Eigentlicher Inhalt beginnt auf Seite 1
\clearpage              % neue Seite für eigentlichen Inhalt

\chapter{Einführung}\label{introduction}

Routenplanungsspiele haben in der Analyse und Optimierung von Verkehrs- und Kommunikationsnetzwerken eine große Bedeutung.
Bei der Modellierung solcher Netzwerke wurden dabei in der Vergangenheit häufig statische Flüsse zu Rate gezogen, welche jedoch nicht abbilden können, wie sich Veränderungen des Flusses im Laufe der Zeit innerhalb des Netzwerks, wie sie beispielsweise in Straßennetzen auftreten, auf das System auswirken.
Daher geht man zu dynamischen Flüssen (engl. flows over time) über, welche den zeitlichen Verlauf der Kantenbelegung erfassen.

In dieser Arbeit unterwirft man dynamische Flüsse dem Modell der deterministischen Warteschlangen (engl. deterministic queuing model):
Sobald Partikel eine Kante durchlaufen möchten, reihen sich diese zunächst in eine Warteschlange der Kante ein, die mit einer gewissen Kapazität abgearbeitet wird, um dann nach einer zusätzlichen Verzögerung am Zielknoten anzukommen.
Dabei kann man sich die Kanten eines Netzwerk als Transportbänder vorstellen, die eine bestimmte Kapazität und eine bestimmte Verzögerung haben.
Die Kapazität entspricht der Breite des Bands, welche die Gütermenge, die das Band pro Zeiteinheit transportieren kann,  beschränkt, und die Verzögerung entspricht der Transportzeit, die das Band benötigt, um Güter vom Startknoten zum Zielknoten zu tragen.
Weiter gibt es im Netzwerk eine Quelle, also einen Knoten, an dem kontinuierlich Partikel entstehen, und eine Senke, in dem Partikel aus dem Netzwerk abfließen.

Diesen Flüssen gibt man nun die folgende spieltheoretische Interpretation:
Jedes infinitesimal kleine Partikel, das in einem Fluss an der Quelle zu einem bestimmten Zeitpunkt entsteht, wird als Spieler in einem nicht-atomaren Auslastungsspiel interpretiert:
Dabei verfolgt jeder Spieler das Ziel, in möglichst kurzer Zeit nach Erscheinen an der Quelle zur Senke des Netzwerks zu gelangen.
Jeder Spieler besitzt dabei bereits zur Zeit seines Entstehens über alle nötigen Informationen, die seine tatsächliche Ankunftszeit an der Senke beeinflussen -- das heißt, er weiß über den künftigen Verlauf der Warteschlangen an den Kanten Bescheid --, und kann diese nutzen um einen zur aktuellen Zeit kürzesten Quelle-Senke-Pfad zu ermitteln.
Da es sich bei den Spielern um infinitesimal kleine Partikel handelt, verändert sich der dynamische Fluss durch die Entscheidung eines einzelnen Partikels nicht.

Ein Nash-Gleichgewicht wird für eine endliche Anzahl an Spielern für gewöhnlich als eine Strategiewahl charakterisiert, in der kein Spieler seine Strategie bei Beibehalten der Wahl der anderen Spieler ändern kann, um seine eigenen Kosten echt zu verringern.
Im Fall dynamischer Flüsse wird dieser Begriff etwas aufgeweicht:
So wird ein dynamischer Fluss ein Nash-Gleichgewicht genannt, falls eine Kante von $v$ nach $w$
fast nur zu Zeitpunkten genutzt wird, zu denen die Kante auf einem kürzesten  Pfad von der Quelle zu $w$ liegt.

Diese Arbeit gibt zunächst eine formale Definition von dynamischen Flüssen in Kapitel~\ref{chapter-dynamic-flows} sowie eine Charakterisierung von Nash-Gleichgewichten in Kapitel~\ref{chapter-nash-flows}.
Danach widmet sie sich der Berechnung ebensolcher Gleichgewichte für den Fall, dass der Netzwerkzufluss an der Quelle konstant ist.
Dazu werden jedoch einige Kenntnisse über zwei Klassen von statischen Flüssen benötigt:
So werden in Kapitel~\ref{chapter-min-con-flows} zunächst Resultate der Analyse von auslastungsminimalen Flüssen zusammengetragen:
Dies sind $b$-Flüsse, die die maximale Auslastung aller Kanten, d.h. das Verhältnis von Fluss und Kapazität einer Kante, minimieren.
Es wird gezeigt, dass das dazu duale Problem ist, einen sogenannten dünnsten Schnitt zu finden, der das Verhältnis vom Nettoangebot und der Kapazität der ausgehenden Kanten des Schnittes maximiert.
Außerdem wird ein Verfahren vorgestellt, mit dem man auslastungsminimale Flüsse in polynomieller Zeit exakt bestimmen kann.
Weiter führt die Analyse der Ableitung dynamischer Nash-Flüssen zum Konzept sogenannter schmaler Flüsse, die in Kapitel~\ref{chapter-thin-flows}.
Dies sind $b$-Flüsse in einem Netzwerk mit einer einzigen Quelle $s$ und sogenannten zurücksetzenden Kanten, in denen alle $s$-$v$-Pfade, die in der Quelle starten und positiven Fluss haben, minimale Auslastung haben.
Nach einer Charakterisierung dieser Flüsse wird ein polynomieller Algorithmus erarbeitet, der schmale Flüsse in Netzwerken ohne zurücksetzende Kanten berechnet.
Außerdem wird diskutiert, ob das allgemeine Problem mit zurücksetzenden Kanten in der Komplexitätsklasse \PPAD\ liegt.
Schließlich leitet man mithilfe der Existenz solcher schmaler Flüsse ein mögliches Verfahren ab, mit dem Nash-Flüsse bei konstantem Netzwerkzufluss berechnet werden sollen.

\chapter{Dynamische Flüsse}\label{sec-dynamic-flows}

Zunächst werden einige grundlegende Begriffe eingeführt:

\begin{definition}[Netzwerk]
	Ein \emph{Netzwerk} $(G, u, s, t, \tau)$ ist ein gerichteter, endlicher Graph $G=(V,E)$ mit Ursprung $s\in V$, also ein Graph in dem $s$ alle Knoten in $V$ erreichen kann.
	Dabei heißen $s$ die \emph{Quelle} und $t\in V$ die \emph{Senke} des Netzwerks.
	Jeder Kante $e\in E$ werden eine Kapazität $u_e > 0$ und eine Verzögerungszeit $\tau_e\geq 0$ zugeordnet, sodass alle Zykel $C$ eine positive Gesamtverzögerung $\sum_{e\in C}\tau_e$ haben.
\end{definition}

\begin{definition}[statischer Fluss]
	Eine Kantenbewertung $f\in\R_{\geq 0}^E$ in einem gerichteten Graphen $(V, E)$ heißt \emph{statischer Fluss}.
	Existieren zwei Knoten $s,t\in V$, sodass
	$$ \sum_{e\in\delta^+(v)} f_e - \sum_{e\in\delta^-(v)} f_e$$ für alle Knoten $v\in V\setminus\{ s, t \}$ verschwindet, für $s$ nicht-negativ und für $t$ nicht-positiv ist, so ist $f$ ein \emph{statischer $s$-$t$-Fluss mit Wert} $v(f):=\sum_{e\in\delta^+(s)} f_e - \sum_{e\in\delta^-(s)}f_e$.
\end{definition}

Man bemerke, dass hier ein statischer Fluss keine Kapazitätsbedingung erfüllen muss.

\begin{definition}
	Der Funktionenraum $\mathfrak{F}_0$ sei die Menge aller Funktionen $g: \R \to \R_{\geq 0}$, die lokal integrierbar bzgl. des Lebesgue-Maßes sind, also $\int_a^b |g(t)| \diff t< \infty$ für beliebige beschränkte Intervalle $(a,b)$ erfüllen, und auf der negativen Achse verschwinden, das heißt, es gilt $g(t)=0$ für $t<0$.
\end{definition}

\begin{definition}[Dynamischer Fluss]
	Ein \emph{dynamischer Fluss $f=(f^+, f^-)$} ist ein Paar zweier über die Kanten $E$ eines Netzwerks indizierten Familien mit $f^+_e,f^-_e\in\mathfrak F_0$ für alle $e\in E$.
	Dabei bezeichnen $f_e^+(\theta)$ bzw. $f_e^-(\theta)$ die \emph{Zu- bzw. Abflussrate an Kante $e\in E$ zum Zeitpunkt $\theta\in\R$}.
	
	Der (kumulative) \emph{Zu- bzw. Abfluss an einer Kante $e$ bis zum Zeitpunkt $\theta$} sei definiert durch $F^+_e(\theta):=\int_0^\theta f^+_e(t) \diff t<\infty$ bzw. $F^-_e(\theta):=\int_0^\theta f^-_e(t) \diff t<\infty$.
	
	Die \emph{(Länge der) Warteschlange $z_e(\theta)$ und die Wartezeit $q_e(\theta)$ an Kante $e\in E$ zum Zeitpunkt $\theta\in\R$} seien gegeben durch $z_e(\theta):= F_e^+(\theta) - F_e^-(\theta + \tau_e)$ und $q_e(\theta) := z_e(\theta) / u_e$.
	
	Man beschreibe die \emph{Austrittszeit $T_e(\theta)$ aus einer Kante $e\in E$ bei Eintrittszeit $\theta$}, zu der ein Partikel eine Kante verlässt, die es zum Zeitpunkt $\theta$ betreten hat, als $T_e(\theta):=\theta + q_e(\theta) + \tau_e$.
\end{definition}

Der Definition der Austrittszeit kann man bereits entnehmen, wie sich Partikel, die an einer Kante $e$ ankommen, verhalten sollen:
Nachdem sie zur Zeit $\theta$ die Kante betreten haben, müssen sie sich zunächst in eine Warteschlange einreihen, welche mit der Kapazität $u_e$ nach dem FIFO-Prinzip abgebaut wird.
Nachdem diese Wartezeit $q_e(\theta)$ vorüber ist, vergeht eine weitere konstante Verzögerungszeit $\tau_e$, bevor sie wieder aus der Kante austreten.
Des Weiteren müssen sich die Partikel bereits sofort bei der Ankunft an einem Knoten entscheiden, in welche Kante sie eintreten wollen, und können nicht an einem Knoten verweilen.

Man kann sich die Kanten also als Transportbänder vorstellen, deren Breite die Transportkapazität bestimmt und deren Länge die Verzögerungszeit darstellt.
Überschreiten die Güter die Kapazität, so bildet sich eine Warteschlange vor dem Band.

Dies führt zu folgender Definition der Zulässigkeit dynamischer Flüsse:

\begin{definition}[Zulässiger dynamischer Fluss]
	Ein dynamischer Fluss $f=(f^+, f^-)$ heißt \emph{zulässig}, falls er die folgenden Eigenschaften erfüllt:
	\begin{enumerate}[label=(F\arabic*)]
		\item\label{def-feasible-flow-capacity} Keine Abflussrate übersteigt die Kapazität, d.h. $\forall e\in E, \theta\in\R: f_e^-(\theta)\leq u_e$.
		\item\label{def-feasible-flow-no-negative-flow} Fluss verlässt eine Kante nur, falls er sie zuvor betreten hat,\\ d.h. $\forall e\in E, \theta\in\R: F_e^+(\theta) \geq F_e^-(\theta + \tau_e).$
		\item\label{def-feasible-flow-no-flow-at-node} Bis auf Quelle und Senke erfüllt jeder Knoten $v$ Flusserhaltung,\\
		d.h. $\forall\theta\in\R: \sum_{e\in\delta^+(v)}f^+_e(\theta) - \sum_{e\in\delta^-(v)} f_e^-(\theta) = 0$.\\
		Für die Senke $t$ muss dieser Wert nicht-positiv und für die Quelle $s$ nicht-negativ sein und heißt für $s$ \emph{der Zufluss $d(\theta)$ in das Netzwerk}.
		\item\label{def-feasible-flow-queue-with-capacity} Warteschlangen werden mit der Kapazität der Kante abgebaut,\\ d.h. $\forall e\in E, \theta\in\R: z_e(\theta) > 0 \implies f_e^-(\theta + \tau_e) = u_e$.
	\end{enumerate}
\end{definition}

Die folgende Proposition beschreibt wichtige Folgerungen über zulässige dynamische Flüsse:

\begin{proposition}\label{prop-feasible-flow}
	Für eine Kante $e\in E$ und einen zulässigen dynamischen Fluss $f$ gilt:
	\begin{enumerate}[label=(\roman*)]
		\item\label{prop-feasible-flow-T-mon-inc-cont} Die Funktion $\theta \mapsto \theta + q_e(\theta)$ ist monoton wachsend und stetig.
		\item\label{prop-feasible-flow-positive-queue} Für alle $\theta\in\R$ ist die Warteschlange $z_e$ auf dem Intervall $(\theta, \theta + q_e(\theta))$ positiv.
		\item\label{prop-feasible-flow-det-outflow} Zu jeder Zeit $\theta\in\R$ gilt $F_e^+(\theta) = F_e^-(T_e(\theta))$.
		\item\label{prop-feasible-flow-queue-delay} Für alle $\theta_1 \leq \theta_2$ mit $\int_{\theta_1}^{\theta_2} f^+_e(t) \diff t = 0$ und $q_e(\theta_2)>0$ gilt $\theta_1 + q_e(\theta_1) = \theta_2 + q_e(\theta_2)$.
	\end{enumerate}
\end{proposition}
\begin{proof}
	In~\ref{prop-feasible-flow-T-mon-inc-cont} folgt die Stetigkeit bereits aus der Stetigkeit von $F_e^+$ und $F_e^-$.
	Um zu zeigen, dass die Funktion monoton wachsend ist, seien $\theta_1 \leq \theta_2$ gegeben.
	Mit der Monotonie von $F_e^+$ und mit $F_e^-(\theta_2 + \tau_e) = F_e^-(\theta_1+\tau_e) + \int_{\theta_1+\tau_e}^{\theta_2+\tau_e} f_e^-(t)\diff t\leq F_e^-(\theta_1 + \tau_e) + (\theta_2 - \theta_1)u_e$ gilt: 
	\[
		\theta_1 + q_e(\theta_1)
		= \theta_1 + \frac{F_e^+(\theta_1) - F_e^-(\theta_1 + \tau_e)}{u_e}
		\leq \theta_1 + \frac{F_e^+(\theta_2) - F_e^-(\theta_1+\tau_e)}{u_e} \leq \theta_2 + q_e(\theta_2).
	\]
	
	Für $\theta'\in (\theta, \theta+q_e(\theta))$ gilt also $\theta' + q_e(\theta') \geq \theta + q_e(\theta)$, womit $q_e(\theta') \geq \theta + q_e(\theta) - \theta' > 0$ gerade Aussage (ii) beweist.
	
	Aussage (iii) folgt dann mit~\ref{def-feasible-flow-queue-with-capacity} und Aussage~(ii), da
	$\int_{\theta}^{\theta + q_e(\theta)}f_e^-(t + \tau_e) \diff t = q_e(\theta)  u_e = z_e(\theta)$ gilt, und damit ist $F_e^-(T_e(\theta)) = F_e^-(\theta+\tau_e) + \int_{\theta+\tau_e}^{\theta+\tau_e+q_e(\theta)}f_e^-(t)\diff t = F_e^+(\theta)$.
	
	Zu Aussage (iv):
	Für alle $\theta'\in [\theta_1, \theta_2]$ gilt $F_e^+(\theta') = F_e^+(\theta_2)$.
	Also ist die Warteschlange $z_e(\theta') = F_e^+(\theta_2) - F_e^-(\theta' + \tau_e) \geq z_e(\theta_2) > 0$ positiv und nach~\ref{def-feasible-flow-queue-with-capacity} gilt $f_e^-(\theta' + \tau_e)=u_e$.
	Die Differenz der Warteschlangen ist
	$z_e(\theta_1)-z_e(\theta_2)=-F^-_e(\theta_1 + \tau_e) + F^-_e(\theta_2 + \tau_e) = (\theta_2 - \theta_1)u_e$, was $q_e(\theta_1) - q_e(\theta_2) = \theta_2 - \theta_1$ impliziert.
\end{proof}
\chapter{Kürzeste Wege}\label{sec-travel-times}

In diesem Abschnitt wird der Begriff der frühesten Ankunftszeit an einem Knoten eingeführt und erörtert, wann eine Kante $vw$ in einem kürzesten $s$-$w$-Pfad liegt.

\begin{definition}
	Für einen dynamischen Fluss $f$ und einen Pfad $P=(e_1,\dots,e_k)$ definiere $l^P(\theta):=T_{e_k}\circ\dots\circ T_{e_1}(\theta)$ den Zeitpunkt, an dem ein Partikel den Endknoten des Pfads erreicht, falls er den Pfad zum Zeitpunkt $\theta$ betritt.
	
	Für einen Knoten $w\in V$ beschreibe $\mathcal{P}_w$ die Menge aller $s$-$w$-Pfade.
	Dann ist die früheste Ankunft eines Partikels, das zur Zeit $\theta$ bei $s$ startet, gegeben durch $l_w(\theta):=\min_{P\in\mathcal{P}_w}l^P(\theta)$.
	Ein Pfad $P\in \mathcal{P}_w$ heißt \emph{kürzester $s$-$w$-Pfad zur Zeit $\theta$}, falls $l_w(\theta)=l^P(\theta)$.
\end{definition}

\begin{proposition}\label{prop-abs-cont-sur}
	Für einen zulässigen Fluss $f$ sind die Funktionen $F_e^+$ und $F_e^-$ für alle $e\in E$ lokal absolut stetig.
	Die Funktionen $T_e$, $l^P$ sowie $l_v$ sind dabei für alle Kanten $e\in E$, Pfade $P$ in G und Knoten $v\in V$ monoton wachsend, lokal absolut stetig und surjektiv.
\end{proposition}
\begin{proof}
	Nach dem Hauptsatz der Differential- und Integralrechnung für das Lebesgue-Inte\-gral (siehe dazu \cite[Satz 4.14]{Elstrodt2011Abs}) ist $G: [a,b] \to \R, x\mapsto \int_a^x g(t) \diff t$ für eine Lebesgue integrierbare Funktion $g: [a,b] \to \R$ absolut stetig.
	Insbesondere sind also $F_e^+$ sowie $F_e^-$ und damit auch $q_e$ und $T_e$ lokal absolut stetig.
	Als Komposition bzw. punktweises Minimum endlich vieler lokal absolut stetiger Funktionen ist auch $l^P$ bzw. $l_v$ für alle Pfade $P$ und Knoten $v$ lokal absolut stetig.
	Nach Proposition~\ref{prop-feasible-flow}~\ref{prop-feasible-flow-T-mon-inc-cont} ist die Monotonie von $T_e$ bereits gegeben, welche auch die Monotonie von $l^P$ und $l_v$ impliziert.
	Wegen $f_e^+, f_e^-\in\mathfrak{F_0}$ gilt $q_e(\theta)=0$ für $\theta\leq 0$, wodurch auch $\lim_{\theta\to-\infty} T_e(\theta) = - \infty$ folgt.
	Mit $T_e(\theta)\geq \theta$ ergibt sich die Surjektivität von $T_e$.
	Daher sind auch $l^P$ und $l_v$ surjektiv.
\end{proof}

Wie im statischen Szenario von kürzesten Pfaden, gilt auch hier die Dreiecksungleichung: 

\begin{lemma}\label{lemma-dreicksungl}
	Für alle Kanten $vw\in E$ gilt in einem zulässigen Fluss 
	$T_{vw}(l_v(\theta)) \geq l_w(\theta)$.
\end{lemma}
\begin{proof}
	Sei ein kürzester $s$-$v$-Pfad $P$ zum Zeitpunkt $\theta$ gegeben.
	Hängt man an $P$ die Kante $vw$ an, erhält man einen $s$-$w$-Pfad, der zur Eintrittszeit $\theta$ die Ankunftszeit $T_{vw}(l_v(\theta))$ liefert.
	Da $l_w(\theta)$ das Minimum über die Ankunftszeit aller $s$-$w$-Pfade ist, gilt die Behauptung.
\end{proof}

\begin{definition}
	Man bezeichne eine Kante $vw\in E$ als \emph{aktiv zum Zeitpunkt $\theta$}, falls sie auf einem zur Zeit $\theta$ kürzesten $s$-$w$-Pfad liegt, d.h. falls $T_{vw}(l_v(\theta)) = l_w(\theta)$ gilt; sonst nennt man sie \emph{inaktiv zum Zeitpunkt $\theta$}.
	Die Menge $\Theta_e$ ist die abgeschlossene Menge aller Zeitpunkte, zu denen die Kante $e$ aktiv ist.
\end{definition}

Man beachte, dass Teilpfade kürzester Pfade im statischen Sinne wieder kürzeste Pfade sind; im dynamischen Sinne gilt dies nicht unbedingt, jedoch aber in folgendem Teilgraph:

\begin{lemma}\label{lemma-shortest-path-using-active-edges}
	Für einen zulässigen Fluss ist der durch die zur Zeit $\theta$ aktiven Kanten induzierte Teilgraph $G_\theta:=(V, E_\theta)$ zu jeder Zeit ein azyklischer Graph, in dem $s$ jeden Knoten $v\in V$ erreichen kann.
\end{lemma}
\begin{proof}
	Angenommen es existiere ein Zyklus $C=(v_1, \dots, v_n)$ mit $v_1=v_n$ und ausschließlich aktiven Kanten.
	Es ist $l^C(\theta) > \theta$, da für Zyklen eine positive Gesamtverzögerung vorausgesetzt ist.
	Aufgrund der Aktivität aller Kanten in $C$ erzeugt $l_{v_1}(\theta) = l^C(l_{v_1}(\theta)) > l_{v_1}(\theta)$ einen Widerspruch.
	
	Für jeden Knoten $w\neq s$ existiert mindestens eine eingehende aktive Kante -- zum Beispiel die letzte Kante eines kürzesten $s$-$w$-Pfades, welcher wiederum existiert, da $w$ von $s$ aus erreichbar ist.
	Daher ist $w$ von $s$ aus in $G_\theta$ erreichbar.
\end{proof}

\begin{proposition}\label{prop-arrival-times-vector}
	Für einen zulässigen Fluss $f$ ist der Vektor $(l_v(\theta))_{v\in V}$ die eindeutige Lösung des Gleichungssystems
	\[ \tilde{l}_w = \begin{cases}
	\theta, & \text{falls } w=s, \\
	\min\limits_{vw\in \delta^-(w)} T_{vw}(\tilde{l}_v), & \text{sonst}.
	\end{cases} \]
\end{proposition}
\begin{proof}
	Offenbar löst $(l_v(\theta))_{v\in V}$ dieses System, da nach Lemma~\ref{lemma-shortest-path-using-active-edges} jeder Knoten $w\neq s$ eine aktive eingehende Kante hat.
	Für eine Lösung $(\tilde{l}_v)_{v\in V}$ des Gleichungssystems, zeige man $l_w(\theta) = \tilde{l}_w$ für jeden Knoten $w\in V$.
	Dabei ist der Teilgraph $G'=(V, E')$ mit
	\[ E' := \{ vw \in E \mid T_{vw}(\tilde{l}_v ) = \tilde{l}_w \} \]
	ein azyklischer Graph, in dem $s$ jeden Knoten $w\in V$ erreichen kann:
	Zyklen können wegen der positiven Gesamtverzögerung nicht entstehen und jeder Knoten $w\neq s$ hat mindestens eine eingehende Kante $vw$ mit $T_{vw}(\tilde{l}_v) = \tilde{l}_w$.
	Daher ist jede Knotenbewertungen $\tilde{l}_w$ bereits durch einen $s$-$w$-Pfad $P$ in $G'$ festgelegt auf $l^P(\theta)\geq l_w(\theta)$.
	Für einen zur Zeit $\theta$ kürzesten $s$-$w$-Pfad $Q$ gilt außerdem $\tilde{l}_w \leq T^Q(\tilde{l}_s) = T^Q(\theta) = l_w(\theta)$.
\end{proof}

Um $(l_v(\theta))_{v\in V}$ für alle $\theta\in\R$ gleichzeitig zu berechnen, kann der Bellman-Ford-Algorithmus auf den Distanzvektor-Funktionen $(l_v)_{v\in V}$ genutzt werden:
Dazu wird in jeder der $n-1$ Iterationen für jede Kante das punktweise Minimum $l_w := \min\{ l_w, T_{vw}\circ l_v \}$ gebildet.
Sind Operationen auf Funktionen nicht möglich oder zu teuer, so kann $(l_v(\theta))_{v\in V}$ für ein spezielles $\theta\in\R$ mittels Dijkstra-Algorithmus ermittelt werden, wobei man die Kosten einer Kante $vw$ erst bei Scanning von $v$ in der Form $q_{vw}(l_v(\theta)) + \tau_{vw}$ berechnet.
\chapter{Dynamische Nash-Flüsse}\label{chapter-nash-flows}

Dieser Abschnitt dient dazu, Nash-Gleichgewichte im Kontext dynamischer Flüsse formal einzuführen.
Dabei hilft die Anschauung, dass Partikel, die zur Zeit $\theta$ an der Quelle erscheinen, in einem Nash-Gleichgewicht möglichst früh -- also zum Zeitpunkt $l_t(\theta)$ -- an der Senke ankommen.

\section{Charakterisierung dynamischer Nash-Flüsse}


Für die Einführung dynamischer Nash-Flüsse benötigt man weitere Definitionen: 

\begin{definition}
	Für eine Kante $vw\in E$ bezeichne $x_{vw}^+(\theta)\coloneq F_{vw}^+(l_v(\theta))$ den Zufluss bis zur frühestmöglichen Ankunftszeit von Partikeln in $v$, die zur Zeit $\theta$ in $s$ starten.\\
	Dagegen bezeichne $x_{vw}^-(\theta)\coloneq F^-_{vw}(l_w(\theta))$ den Abfluss bis zur frühestmöglichen Ankunftszeit von Partikeln in $w$, die zur Zeit $\theta$ in $s$ starten.
	
	Für einen Knoten $v\in V$ sei $b_v(\theta)\coloneq \sum_{e\in\delta^+(v)} x_e^+(\theta) - \sum_{e\in\delta^-(v)} x_e^-(\theta)$ die Balance des Knotens $v$ zum Zeitpunkt $\theta$.
\end{definition}

\begin{remark}\label{remark-x^-leqx^+}
	In einem zulässigen Fluss gilt nach Proposition~\ref{prop-feasible-flow}~\ref{prop-feasible-flow-det-outflow} und mit der Monotonie von $F_{vw}^-$ bereits 
	\[
	x_{vw}^-(\theta) = F_{vw}^-(l_w(\theta)) \leq F_{vw}^-(T_{vw}(l_v(\theta)))=F_{vw}^+(l_v(\theta)) = x_{vw}^+(\theta).
	\]
\end{remark}

\begin{lemma}\label{lemma-balance-0}
	Für einen zulässigen dynamischen Fluss $f$ gilt $b_v(\theta)=0$ für alle Knoten $v\in V\setminus\{ s,t \}$ und alle $\theta\in\R$.
\end{lemma}
\begin{proof}
	Unter Benutzung von~\ref{def-feasible-flow-no-flow-at-node} folgere man für alle $v\in V\setminus \{ s, t\}$ und $\theta\in\R$:
	\[ \sum_{e\in\delta^-(v)} x_e^-(\theta) = \int_{0}^{l_v(\theta)} \sum_{e\in\delta^-(v)} f_e^-(t) \diff t = \int_{0}^{l_v(\theta)} \sum_{e\in\delta^+(v)} f_e^+(t) \diff t = \sum_{e\in\delta^+(v)}x_e^+(\theta). \]
\end{proof}

\begin{notation}
	 $M^c\coloneq \R\setminus M$ bezeichne das Komplement einer Menge $M\subseteq\R$ und $\overline{M}$ bezeichne den Abschluss von $M$.
\end{notation}

\begin{definition}\label{def-flow-along-active-edges}
	Man sage, der Fluss $f$ \emph{fließe nur entlang aktiver Kanten}, falls $f_{vw}^+$ fast überall auf $l_v(\Theta_{vw}^c)$ für alle Kanten $vw\in E$ verschwindet.
\end{definition}

\begin{remark}
	Diese Definition weicht von der Definition von Koch und Skutella ab und entspricht derjenigen aus~\cite[Definition 1]{Cominetti2015}:
	Nach \cite[Definition 2]{Koch2011} sagt man, $f$ \emph{sende Fluss nur entlang aktuell kürzester Pfade}, falls $f_{vw}^+\circ l_v$ fast überall auf $\Theta_{vw}^c$ verschwindet für alle Kanten $vw$.

	Entspricht $f$ dieser Definition, so auch Definition~\ref{def-flow-along-active-edges}: 
	Da $l_v$ nach Proposi\-tion~\ref{prop-abs-cont-sur} absolut stetig ist, bildet es nach~\cite[Kap. VII, Aufgabe 4.9]{Elstrodt2011} Nullmengen wieder auf Nullmengen ab, weshalb folgende Menge eine Nullmenge ist: \[ l_v(\{ \theta \in \Theta_{vw}^c \mid f_{vw}^+ (l_v(\theta)) > 0 \}) = \{ \xi \in l_v(\Theta_{vw}^c) \mid f_{vw}^+ (\xi) > 0 \}. \]
	 
	Koch und Skutella zeigen im Beweis von~\cite[Lemma 1]{Koch2011} die entsprechende Äquivalenz von Lemma~\ref{lemma-only-active-edges} (i) und (iii) -- jedoch in (i) unter Verwendung ihrer Definition -- und
	verwenden bei der Implikation (iii)$\Rightarrow$(i) das Argument, dass für jede Kante $vw\in E$ und alle $\theta\in \Theta_{vw}^c$ eine Umgebung $U$ von $\theta$ existiert, sodass $f_{vw}^+$ fast überall in $l_v(U)$ verschwindet.
	Dies reicht aber nicht aus, um zu zeigen, dass $f_{vw}^+(l_v(\theta))=0$ für fast alle $\theta\in\Theta_{vw}^c$ gilt:
	So kann $f_{vw}^+(l_v(\theta))$ für ein $\theta\in\Theta_{vw}^c$ positiv sein und $l_v$ konstant in einer Umgebung um $\theta$.
	Dann ist $f_{vw}^+ \circ l_v$ in einer Umgebung um $\theta$ positiv, was im Widerspruch zur Forderung steht.
	
	Dies wurde in~\cite[Example 2]{Cominetti2015} dazu genutzt, einen Beispielfluss anzugeben, der beweist, dass die Forderung von Koch und Skutella sogar echt stärker ist.
\end{remark}

Für eine äquivalente Umschreibung dieser Definition benötigen wir folgendes Lemma der Maßtheorie:

\begin{lemma}\label{lemma-vanishes-intervals}
	Seien $g: \R \to \R_{\geq 0}$ eine lokal Lebesgue-integrierbare Funktion und $((a_i, b_i))_{i\in I}$ eine Familie offener Intervalle.
	Dann verschwindet $g$ fast überall auf der Vereinigung $\Theta\coloneq \bigcup_{i\in I} (a_i, b_i)$ genau dann, wenn es für alle $i\in I$ fast überall auf $(a_i, b_i)$ verschwindet.
\end{lemma}
\begin{proof}
	Verschwindet $g$ fast überall auf $\Theta$, so erst recht auf jedem Intervall $(a_i, b_i)$.
	Es verschwinde $g$ also auf allen Intervallen $(a_i, b_i)$.
	Die Funktion $\mu(A)\coloneq \int_A g \diff \lambda$ definiert ein Maß auf den Borelmengen~$\mathfrak{B}$.
	Da jede offene Menge $O\subseteq\R$ bereits $\sigma$-kompakt ist, also eine Darstellung als abzählbare Vereinigung kompakter Mengen -- hier $O=\bigcup_{n\in\N}(\{ x \in\R \mid d(x, O^c) \geq 1/n \} \cap [-n, n] )$ -- besitzt, ist jede offene Menge nach~\cite[Kap. VIII, 1.2 Folgerungen (e)]{Elstrodt2011} innen regulär.
	Das heißt, es gilt
	\[ \mu(O)=\sup\{ \mu(K) \mid K\subseteq O \text{ kompakt} \} \]
	für offene Mengen $O\subseteq\R$.
	Für ein kompaktes $K\subseteq \Theta$ existiert eine endliche Teil\-über\-deckung $\bigcup_{i=1}^n (a_i, b_i) \supseteq K$, für die $\mu(K) \leq \sum_{i=1}^{n} \mu((a_i, b_i)) = \sum_{k=1}^{n} \int_{a_i}^{b_i} g(t) \diff t = 0$ gilt.
	Also ist auch $\mu(\Theta)=0$.
\end{proof}

\begin{lemma}\label{lemma-only-active-edges}
	Für einen zulässigen Fluss $f$ sind folgende Aussagen äquivalent:
	\begin{enumerate}[label=(\roman*)]
		\item Der Fluss $f$ fließt nur entlang aktiver Kanten.
		\item Für jede Kante $vw\in E$ und für fast alle $\xi\in\R$ mit	$f_{vw}^+(\xi)>0$ gilt $\xi \in l_v(\Theta_{vw})$.
		\item Für jede Kante $e\in E$ und für alle $\theta\in\R$ gilt $x_e^+(\theta) = x_e^-(\theta)$.
	\end{enumerate}
\end{lemma}
\begin{proof}
	(i)$\Leftrightarrow$(ii): Bedingung~(ii) gilt genau dann, wenn $f_{vw}^+$ fast überall auf $l_v(\Theta_{vw})^c$ verschwindet.
	Daher genügt es, zu zeigen, dass sich $l_v(\Theta_{vw})^c$ und $l_v(\Theta_{vw}^c)$ nur um eine Nullmenge voneinander unterscheiden.
	Mit der Surjektivität von $l_v$ gilt $l_v(\Theta_{vw})^c\subseteq l_v(\Theta_{vw}^c)$.
	
	Des Weiteren ist $S\coloneq l_v(\Theta_{vw}^c)\setminus l_v(\Theta_{vw})^c = l_v(\Theta_{vw}^c)\cap l_v(\Theta_{vw})$ eine Teilmenge von $l_v(\Q)$:
	Für ein $\xi\in S$ gibt es $\theta\in\Theta_{vw}^c$ und $\theta'\in\Theta_{vw}$ mit $l_v(\theta)=\xi=l_v(\theta')$.
	Da $\theta\neq\theta'$ ist, existiert ein $\theta_q\in\Q\cap(\theta,\theta')$.
	Wegen der Monotonie von $l_v$ gilt $l_v(\theta_q)=\xi$, womit $\xi\in l_v(\Q)$ folgt.
	Also unterscheiden sich die beiden Mengen nur um eine abzählbare Menge.
	
	(i)$\Leftrightarrow$(iii): Sei eine Kante $vw\in E$ gegeben.
	Für ein $\theta\in\R$ bezeichne $\omega_\theta\leq \theta$ den spätesten Startzeitpunkt, sodass man unter Benutzung von $vw$ zum Zeitpunkt $l_w(\theta)$ zu $w$ gelangt:
	\[ \omega_\theta\coloneq \max\{ \omega\leq\theta \mid l_w(\theta) = T_{vw}(l_v(\omega)) \}. \]
	
	Es gilt $\Theta_{vw}^c = \bigcup_{\theta\in\R} (\omega_\theta, \theta)$:
	Für $\theta\in\Theta_{vw}^c$ gilt $T_{vw}(l_v(\theta)) > l_w(\theta)$.
	Aufgrund der Stetigkeit von $T_{vw}\circ l_v$ und $l_w$ existiert ein $\varepsilon>0$, das $T_{vw}(l_v(\theta')) > l_w(\theta+\varepsilon)$ für alle $\theta'\in[\theta,\theta+\varepsilon]$ erfüllt.
	Also liegt $\theta$ in $(\omega_{\theta+\varepsilon}, \theta+\varepsilon)$.
	Ist umgekehrt $\theta'\in (\omega_\theta,\theta)$, so gilt aufgrund der Monotonie $T_{vw}(l_v(\theta'))\geq T_{vw}(l_v(\omega_\theta)) = l_w(\theta)\geq l_w(\theta')$.
	Die erste Ungleichung kann nicht mit Gleichheit erfüllt sein, da $\omega_\theta$ maximal mit der Eigenschaft $T_{vw}(l_v(\omega)) = l_w(\theta)$ ist, wodurch $\theta'\in\Theta_{vw}^c$ folgt.
	
	Mit $l_v(\Theta_{vw}^c) = \bigcup_{\theta\in\R}(l_v(\omega_\theta),l_v(\theta))$ verschwindet $f_{vw}^+$ nach Lemma~\ref{lemma-vanishes-intervals} genau dann fast überall auf $l_v(\Theta_{vw}^c)$, wenn es für alle $\theta\in\R$ fast überall auf $(l_v(\omega_\theta),l_v(\theta))$ verschwindet.
	Nach Proposition~\ref{prop-feasible-flow}~\ref{prop-feasible-flow-det-outflow} gilt dies wiederum genau dann, wenn 	$F_{vw}^+(l_v(\theta))-F_{vw}^-(l_w(\theta))$ für alle $\theta\in\R$ verschwindet.
\end{proof}

\begin{definition}
	Man sage, ein zulässiger dynamischer Fluss $f$ \emph{fließe ohne Über\-holungs\-möglichkeiten}, falls $b_s(\theta) = -b_t(\theta)$ für alle $\theta\in\R$ gilt.
\end{definition}

Dabei betrachte man folgende Intuition:
 Partikel, die zur Zeit $\theta\in\R$ bei $s$ starten und sich auf einem kürzesten Weg zu $t$ bewegen, also zur Zeit $l_t(\theta)$ in $t$ ankommen, überholen andere Partikel, falls $b_s(\theta) > -b_t(\theta)$ gilt.
Falls jedoch $b_s(\theta) < - b_t(\theta)$ gilt, wurde das Partikel bereits von anderen überholt.
Ein Nash-Gleichgewicht sollte diese Eigenschaft daher erfüllen.

\newcommand{\newv}{\mathbf{v}}
\begin{lemma}\label{lemma-b-graph}
	Seien ein dynamischer Fluss $f$ in einem Graphen $G=(V,E)$ und ein Zeitpunkt $\theta\in\R$ gegeben.
	Der Graph $H$ entstehe aus $G$, indem man jede Kante $vw\in E$ aus $G$ durch einen neuen Knoten $\newv_{vw}$ und zwei Kanten $v\newv_{vw}$ und $\newv_{vw}w$ ersetze.
	Der statische Fluss $g$ auf $H$ sei definiert durch
	\[ g_{v\newv_{vw}} \coloneq x_{vw}^+(\theta) \text{ und } g_{\newv_{vw}w} \coloneq x_{vw}^-(\theta) \text{ für alle $vw\in E$} \]
	und der Balancevektor $b$ auf $H$ sei gegeben durch $b_v\coloneq b_v(\theta)$ für alle $v\in V$ und $b_{\newv_e}\coloneq x_e^-(\theta) - x_e^+(\theta)$ für alle $e\in E$.
	Dann gelten die folgenden Aussagen:
	
	\begin{enumerate}[label=(\roman*)]
		\item Der Fluss $g$ ist ein statischer $b$-Fluss.
		\item\label{lemma-b-graph-imp} Ist $f$ zulässig, so gilt $\forall e\in E : x_e^+(\theta) = x_e^-(\theta)\iff b_s(\theta) + b_t(\theta) = 0$.
	\end{enumerate}
\end{lemma} 
\begin{proof}
	Zu (i): Um zu zeigen, dass die Summe über die Balanceeinträge verschwindet, erkenne man, dass der Anteil einer Kante $e\in E$ in $b(V)$ gerade $x_e^+(\theta) - x_e^-(\theta)$ ist.
	Damit gilt:
		\[ \sum_{v\in V}b_v + \sum_{e\in E} b_{\newv_e} = \sum_{e\in E} (x_e^+(\theta) - x_e^-(\theta) + x_e^-(\theta) - x_e^+(\theta)) = 0. \]
		Es bleibt zu zeigen, dass $g$ bezüglich $b$ Flusserhaltung gewährt.
		Für die Knoten der Form $\newv_{vw}$ gilt dies, da $g_{\newv_{vw}w} - g_{v\newv_{vw}} = x_{vw}^-(\theta) - x_{vw}^+(\theta) = b_{\newv_{vw}}$.
		Für $v\in V$ gilt nach Konstruktion
		\[ b_v =
		\sum_{e\in\delta^+_G(v)} x_{e}^+(\theta) - \sum_{e\in\delta^-_G(v)} x_{e}^-(\theta) =
	\sum_{e\in\delta_H^+(v)} g_e - \sum_{e\in\delta^-_H(v)}g_e
		. \]
	
	Zu (ii): Tatsächlich benötigt man aus (i) nur die Eigenschaft, dass die Summe über die Einträge des Balancevektors verschwindet.
	Mit Lemma~\ref{lemma-balance-0} gilt wegen der Zulässigkeit von $f$ sogar $b_s(\theta)+b_t(\theta) + \sum_{e\in E} b_{\newv_e} = 0$.
	
	Angenommen, es gelte $x_e^+(\theta) = x_e^-(\theta)$ für alle $e\in E$.
	Dann sind auch alle $b_{\newv_e} = 0$ und es gilt $b_s(\theta) + b_t(\theta) = 0$.
	Setzt man $b_s(\theta) + b_t(\theta) = 0$ voraus, so ist $\sum_{e\in E} b_{\newv_e} = 0$ und, da $f$ zulässig ist, gilt $x_e^-(\theta)\leq x_e^+(\theta)$ nach Bemerkung~\ref{remark-x^-leqx^+}.
	Daher gilt $b_{\newv_e}\leq 0$ für alle $e\in E$, weshalb $b_{\newv_e}$ für alle $e\in E$ verschwinden muss.
\end{proof}

Die Ergebnisse aus Lemma~\ref{lemma-only-active-edges} und Lemma~\ref{lemma-b-graph} werden im folgenden Theorem gesammelt, welches Nash-Gleichgewichte in dynamischen Flüssen charakterisiert:

\begin{theorem}[Charakterisierung dynamischer Nash-Flüsse]\label{thm-equivalencies-nash-flow}
	Ist $f$ ein zu\-läs\-siger dynamischer Fluss, so sind die folgenden Aussagen äquivalent:
	\begin{enumerate}[label=(\roman*)]
		\item Der Fluss $f$ fließt nur entlang aktiver Kanten.
		\item Für alle Kanten $e\in E$ und zu jeder Zeit $\theta\in\R$ gilt $x_e^+(\theta) = x_e^-(\theta)$.
		\item Der Fluss $f$ fließt ohne Überholungsmöglichkeiten.
	\end{enumerate}
	Gilt eine dieser Aussagen, so nennt man $f$ einen \emph{dynamischen Nash-Fluss}.
\end{theorem}

\section{Eigenschaften dynamischer Nash-Flüsse}

In diesem Abschnitt werden einige Ergebnisse über Nash-Flüsse gesammelt, die in Abschnitt~\ref{sec-nash-flow-extension} benötigt werden.

\begin{remark}\label{remark-s-t-flow}
	Nach Lemma~\ref{lemma-balance-0} ist der statische Fluss $x(\theta)$ für einen Nash-Fluss $f$ mit $x_e(\theta)\coloneq x_e^+(\theta)=x_e^-(\theta)$ für alle $\theta\in\R$ ein statischer $s$-$t$-Fluss.
	Wegen der Monotonie von $x_e$ ist auch $x(\theta_2) - x(\theta_1)$ für $\theta_1 \leq \theta_2$ ein statischer $s$-$t$-Fluss genauso wie $x'(\theta)$, falls $x_e$ für alle $e\in E$ differenzierbar in $\theta$ ist, da Differenzieren die Flusserhaltung erhält und $x_e$ monoton wachsend für alle $e\in E$ ist.
\end{remark}

\begin{lemma}\label{lemma-x-locally-constant}
In einem dynamischen Nash-Fluss ist $x_e$ eingeschränkt auf $\overline{\Theta_e^c}$, also auf den Abschluss der Menge der inaktiven Zeitpunkte von $e$, für jede Kante $e\in E$ lokal konstant.
\end{lemma}
\begin{proof}
Da $\Theta_{vw}^c$ eine in $\R$ offene Menge ist, hat sie eine Darstellung als abzählbare Vereinigung paarweise disjunkter offener Intervalle.
Innerhalb eines solchen Intervalls $(\theta_1, \theta_2)$ gilt $x_{vw}(\theta_2) - x_{vw}(\theta_1) = \int_{l_v(\theta_1)}^{l_v(\theta_2)} f_{vw}^+(t) \diff t = 0$, da $f$ nur entlang aktiver Kanten fließt.
Der Rest folgt mit der Monotonie und Stetigkeit von $x_{vw}$.
\end{proof}

\begin{lemma}\label{lemma-nash-flow-waiting-queue-implies-active-edge}
	Seien ein dynamischer Nash-Fluss $f$, eine Kante $vw\in E$ und ein Zeitpunkt $\theta\in\R$ gegeben.
	Gilt eine der folgenden Aussagen, so ist $vw$ zum Zeitpunkt $\theta$ aktiv:
	\begin{enumerate}[label=(\roman*)]
		\item Die Ableitung $x_{vw}'(\theta)$ existiert und ist positiv.
		\item Die Wartezeit $q_{vw}$ an der Kante $vw$ ist zur Zeit $l_v(\theta)$ positiv.
	\end{enumerate}
	Insbesondere verschwindet die Wartezeit $q_{vw}(l_v(\theta))$ für alle $\theta\in\overline{\Theta_{vw}^c}$.
\end{lemma}
\begin{proof}
	Zu Aussage (i): Angenommen, $vw$ wäre zum Zeitpunkt $\theta$ nicht aktiv, so würde wegen der Offenheit von $\Theta_{vw}^c$ und Lemma~\ref{lemma-x-locally-constant} die Ableitung $x_{vw}'(\theta)$ verschwinden.
	
	Für Aussage (ii) zeige man $T_{vw}(l_v(\theta)) \leq l_w(\theta)$.
	Sei $\theta_1$ der früheste Zeitpunkt mit $x_{vw}^+(\theta_1)= x_{vw}^+(\theta)$.
	Dieser existiert, da $l_v$ nach Proposition~\ref{prop-abs-cont-sur} surjektiv ist.
	Dann ist $\theta_1\in \Theta_{vw}$ nach Lemma~\ref{lemma-x-locally-constant}.
	Außerdem ist $\theta_1 \leq \theta$ wegen der Monotonie von $F_{vw}^+ \circ l_v$.
	Nach Aussage (i) gilt nun $T_{vw}(l_v(\theta_1)) = l_w(\theta_1)$.
	Nach Proposition~\ref{prop-feasible-flow}~\ref{prop-feasible-flow-queue-delay} ist $T_{vw}(l_v(\theta_1)) = T_{vw}(l_v(\theta))$ und mit der Monotonie von $l_w$ folgt $T_{vw}(l_v(\theta))\leq l_w(\theta)$.
\end{proof}

\begin{proposition}\label{prop-nash-flow-s-t-path-decomposable}
	Für einen dynamischen Nash-Fluss $f$ und zwei Zeitpunkte $\theta_1,\theta_2$ mit $\theta_1 \leq \theta_2$ ist der statische $s$-$t$-Fluss $x(\theta_2) - x(\theta_1)$ eine Komposition von $s$-$t$-Wegen.
\end{proposition}
\begin{proof}
	Sei $\theta$ das Infimum aller Zeitpunkte $\xi\geq\theta_1$, zu denen $x(\xi) - x(\theta_1)$ nicht in $s$-$t$-Wege zerlegbar ist.
	Man nehme $\theta \leq \theta_2$ an.
	Da inaktive Kanten zum Zeitpunkt $\theta$ bereits kurz vor und noch kurz nach $\theta$ inaktiv sind, existiert ein Intervall $[\theta - \varepsilon, \theta + \varepsilon]$, in dem keine inaktive Kante aktiv wird.
	Außerdem existiert für $\xi_0 \coloneq \max \{ \theta_1, \theta - \varepsilon \}$ eine $s$-$t$-Wegezerlegung von $x(\xi_0) - x(\theta_1)$.
	
	Für einen Pfad $P$ und einen statischen Fluss $g$ sei $g^P \coloneq \min_{e\in P} g_e$ der Fluss, der auf dem Pfad $P$ fließt.
	Für einen Zyklus $C$ ist $(x(\xi) - x(\xi_0))^C = 0$ für $\xi\in [\xi_0, \theta+\varepsilon]$, da aufgrund der Azyklizität von $G_{\xi_0}$ eine Kante $e\in C$ des Zyklus existiert, die zur Zeit $\xi_0$ und damit in ganz $[\xi_0, \theta+\varepsilon]$ inaktiv ist, wodurch $x_e(\xi_0) = x_e(\xi)$ nach Lemma~\ref{lemma-x-locally-constant} folgt.
	
	Also hat der $s$-$t$-Fluss $x(\xi) - x(\xi_0)$ für $\xi\in [\xi_0, \theta + \varepsilon]$ keinen Zyklus mit positivem Fluss und besitzt daher eine $s$-$t$-Wegezerlegung.
	Addiert man diese zur $s$-$t$-Wegezerlegung von $x(\xi_0) - x(\theta_1)$, so erhält man eine $s$-$t$-Wegezerlegung von $x(\xi) - x(\theta_1)$, was für $\xi > \theta$ einen Widerspruch zur Definition von $\theta$ darstellt.
\end{proof}

\begin{corollary}
	Für einen dynamischen Nash-Fluss $f$ ist der statische $s$-$t$-Fluss $x(\theta)$ zu jeder Zeit $\theta$ eine Komposition von $s$-$t$-Wegen.
\end{corollary}
\begin{proof}
	Nach Proposition~\ref{prop-abs-cont-sur} existiert ein Zeitpunkt $\xi_0$ mit $l_v(\xi_0) \leq 0$ für alle Knoten $v\in V$.
	Für $\theta \leq \xi_0$ ist $x(\theta)$ der Nullfluss und offenbar in $s$-$t$-Wege zerlegbar, da die Funktionen $f_e^+$ und $f_e^-$ links der $y$-Achse verschwinden.
	Ist $\theta > \xi_0$,  so ist $x(\theta)= x(\theta) - x(\xi_0)$ nach Proposition~\ref{prop-nash-flow-s-t-path-decomposable} in $s$-$t$-Wege zerlegbar.
\end{proof}
\chapter{Schmale Flüsse mit Zurücksetzen}\label{sec-thin-flows}
\newcommand*{\PlH}{\makebox[1ex]{\textbf{$\cdot$}}}

\section{Definition und Eigenschaften}

Der Abschnitt beginnt mit der Einführung einer neuen Klasse statischer Flüsse:

\begin{definition}[$s$-Fluss]
	Ein $b$-Fluss $f$ in einem Graphen $(V, E)$ heißt \emph{$s$-Fluss} für ein $s\in V$, falls $b_v\leq 0$ für alle $v\in V\setminus \{ s \}$ gilt.
	Dabei bezeichnet $b_s$ den \emph{Wert von $f$}. 
\end{definition}

\begin{definition}[Netzwerk mit Zurücksetzen]
	Ein \emph{Netzwerk mit Zurücksetzen auf $E_1$ und Ursprung $s$} ist ein Tupel $(V, E, u, s, E_1)$, wobei $(V, E, u)$ ein azyklisches Netzwerk ist, in dem der Knoten $s\in V$ jeden anderen Knoten erreicht, und $E_1\subseteq E$ die Menge der \emph{zurücksetzenden Kanten} ist.
\end{definition}

\begin{definition}[Auslastung]
	Sei ein $s$-Fluss $f$ in einem Netzwerk mit Zurücksetzen auf $E_1$ und Ursprung $s$ gegeben.
	Dann ist die \emph{Auslastung einer Kante} gegeben durch $f_e/u_e$.
	Die \emph{Auslastungsübertragung einer Kante $vw$} ist gegeben durch \[ \rho_{vw}(l_v, f_{vw}) := \begin{cases}
		\max\{ l_v, f_{vw} / u_{vw} \}, & \text{falls $vw\notin E_1$,}\\
		f_{vw} / u_{vw}, & \text{falls $vw\in E_1$.}
	\end{cases}
	\]
	Die \emph{Auslastung eines $s$-$v$-Pfades $P=(e_1,\dots,e_k)$} ist dann gegeben durch die Verkettung $(\rho_{e_k}(\PlH, f_{e_k})\circ \dots \circ \rho_{e_1}(\PlH, f_{e_1}))(0)$.
	Die zu $f$ \emph{zugehörige Knotenauslastung $l\in\R_{\geq 0}^V$} sei dann für einen Knoten $v\in V$ festgelegt durch die minimale Auslastung eines $s$-$v$-Pfades.
\end{definition}

Betrachtet man einen $s$-$v$-Pfad $P=(e_1, \dots, e_k)$, der Kanten aus $E_1$ enthält, so ist die Auslastung von $P$ gerade $\max_{j=i}^k f_{e_j}/u_{e_j}$, wobei $e_i$ die letzte Kante des Pfades in $E_1$ ist.
Daher nennt man die Kanten in $E_1$ die zurücksetzenden Kanten, da sie die Auslastung eines vorangegangen Pfades auf ihre eigene zurücksetzen.

\begin{proposition}\label{prop-congestion-labels-dijkstra}
	Die Knotenauslastungen eines $s$-Flusses $f$ sind gegeben durch die eindeutige Lösung $(\tilde{l}_v)_{v\in V}$ des Gleichungssystems
	\[
	\tilde{l}_w = \begin{cases}
		0, & \text{falls $w=s$,}\\
		\min_{vw\in \delta^-(w)} \rho_{vw}(\tilde{l}_v, f_{vw}), & \text{sonst.}
	\end{cases}
	\]
\end{proposition}
\begin{proof}
	Die Existenz einer Lösung folgt aus der Azyklizität des Netzwerks und wegen der Erreichbarkeit jedes Knotens von $s$ aus.
	
	Sei also $\tilde{l}\in\R_{\geq 0}^V$ die Lösung des Gleichungssystems.
	Der Teilgraph $(V, E')$ mit \[
	E':= \left\{ vw\in E ~\middle\vert~ \rho_{vw}(\tilde{l}_v, f_{vw}) = \min_{uw\in\delta^-(w)}  \rho_{uw}(\tilde{l}_u, f_{uw}) \right\}
	\]
	ebenfalls azyklisch und jeder Knoten ist von $s$ aus erreichbar.
	Man zeige $l_w = \tilde{l}_w$ durch eine Induktion über die Distanz von $w$ zu $s$ bezüglich der Anzahl an Kanten.
	Für $w=s$ gilt offenbar $l_s = 0$.
	Sei nun ein $s$-$w$-Pfad $P$ gegeben, seien $vw$ die letzte Kante und $Q$ das restliche Anfangsstück dieses Pfades.
	Nach Induktionsvoraussetzung ist die Auslastung von $Q$ mindestens $\tilde{l}_v$.
	Daher ist die Auslastung von $P$ mindestens $\rho_{vw}(\tilde{l}_v, f_{vw})$ aufgrund der Monotonie von $\rho_{vw}(\PlH, f_{vw})$, wodurch $l_v \geq \tilde{l}_v$ folgt.
	Außerdem hat ein $s$-$v$-Pfad $P$, der in $(V, E')$ verläuft, hat Auslastung $\tilde{l}_v$.
	Da solch ein Pfad existiert, gilt also $l_v \leq \tilde{l}_v$.
\end{proof}

\newcommand{\problemThinFlow}{\textsc{ThinFlow}}
\begin{definition}[Schmaler Fluss]
	Seien ein $s$-Fluss $f$ in einem Netzwerk mit Zurück\-setzen auf $E_1$ und Ursprung $s$ sowie die durch $f$ induzierten Knotenauslastungen $l$ gegeben.
	Der Fluss $f$ heißt \emph{schmaler Fluss mit Zurücksetzen auf $E_1$}, falls die Auslastung jedes $s$-$v$-Pfades mit positivem Fluss $l_v$ beträgt.
	
	Dabei nennt man das Tupel $(V, E, u, s, b, E_1)$ eine \emph{Instanz des $\problemThinFlow$ Problems}.
\end{definition}

Insbesondere ist bei einem schmalen Fluss die Knotenauslastung $l_v$ eines Knotens $v$ gerade die Auslastung eines jeden $s$-$v$-Pfades mit positivem Fluss und $0$, falls kein solcher existiert.

\begin{remark}
	Ronald Koch betrachtet in~\cite{Koch2012} allgemeinere schmale Flüsse: Dort ist das zugrundeliegende Netzwerk nicht als azyklisch vorausgesetzt.
	Diese Einschränkung wird hier im Sinne von~\cite{Cominetti2015} beibehalten.
\end{remark}

\begin{figure}
	\centering
	\newcommand{\newnode}[4]{\node[lul] (#1) at #2 {#3 \\[0.5em] #4};}
	\begin{subfigure}{\textwidth}
		\begin{tikzpicture}[lul/.style={draw,
			ellipse,
			align=center,
			inner sep=0pt,
			outer sep=4pt,
			text width=7mm,
			minimum height=1.5cm
		}]
		
		\newnode{0}{(0,2)}{$a$}{$32$}
		\newnode{1}{(3,2)}{$b$}{$0$}
		\newnode{2}{(5,0)}{$c$}{$-2$}
		\newnode{3}{(7,2)}{$d$}{$-12$}
		\newnode{4}{(10,3)}{$e$}{$-2$}
		\newnode{5}{(10,1)}{$f$}{$0$}
		\newnode{6}{(13,2)}{$g$}{$-16$}
		
		\begin{scope}[-Latex]
		\path [-Latex] (0) edge node[above] {$4$} (1);
		\path [-Latex] (0) edge[bend right] node[above] {$4$} (2);
		\path [-Latex] (0) edge[bend left] node[above] {$4$} (4);
		\path [-Latex] (1) edge node[right] {$2$} (2);
		\path [-Latex] (1) edge node[above] {$2$} (3);
		\path [-Latex] (2) edge node[left] {$2$} (3);
		\path [-Latex] (3) edge node[above] {$1$} (4);
		\path [-Latex] (3) edge node[above] {$2$} (5);
		\path [-Latex] (4) edge node[above] {$1$} (6);
		\path [-Latex] (5) edge node[above] {$1$} (6);
		\end{scope}
		\end{tikzpicture}
		\subcaption{Der Ursprungsgraph. Die Kanten $e$ sind mit Kapazitäten $u_e$ und die Knoten $v$ mit einem Bezeichner $v$ und einer Balance $b_v$ beschriftet.}
	\end{subfigure}
	\par\bigskip
	\begin{subfigure}{\textwidth}
		\begin{tikzpicture}[lul/.style={draw,
			ellipse,
			align=center,
			inner sep=0pt,
			outer sep=4pt,
			text width=7mm,
			minimum height=1.5cm
		}]
		
		\newnode{0}{(0,2)}{$a$}{$0$}
		\newnode{1}{(3,2)}{$b$}{$2{,}75$}
		\newnode{2}{(5,0)}{$c$}{$2{,}75$}
		\newnode{3}{(7,2)}{$d$}{$5$}
		\newnode{4}{(10,3)}{$e$}{$2{,}5$}
		\newnode{5}{(10,1)}{$f$}{$5$}
		\newnode{6}{(13,2)}{$g$}{$8$}
		
		\begin{scope}[-Latex]
		\path [-Latex] (0) edge node[above] {$11/4$} (1);
		\path [-Latex] (0) edge[bend right] node[above] {$11/4$} (2);
		\path [-Latex] (0) edge[bend left] node[above] {$10/4$} (4);
		\path [-Latex] (1) edge node[right] {$1/2$} (2);
		\path [-Latex] (1) edge node[above] {$10/2$} (3);
		\path [-Latex] (2) edge node[left] {$10/2$} (3);
		\path [-Latex] (3) edge node[above] {$0/1$} (4);
		\path [-Latex] (3) edge node[above] {$8/2$} (5);
		\path [-Latex] (4) edge node[above] {$8/1$} (6);
		\path [-Latex] (5) edge node[above] {$8/1$} (6);
		\end{scope}
		\end{tikzpicture}
		\subcaption{Ein schmaler Fluss $f$ ohne Zurücksetzen.}
		\label{subfig-thin-flow-without-resetting}
	\end{subfigure}
	\par\bigskip
	\begin{subfigure}{\textwidth}
		\begin{tikzpicture}[lul/.style={draw,
			ellipse,
			align=center,
			inner sep=0pt,
			outer sep=4pt,
			text width=7mm,
			minimum height=1.5cm
		},
		scale=1]
		
		\newnode{0}{(0,2)}{$a$}{$0$}
		\newnode{1}{(3,2)}{$b$}{$3$}
		\newnode{2}{(5,0)}{$c$}{$3$}
		\newnode{3}{(7,2)}{$d$}{$5{,}5$}
		\newnode{4}{(10,3)}{$e$}{$2$}
		\newnode{5}{(10,1)}{$f$}{$5{,}5$}
		\newnode{6}{(13,2)}{$g$}{$8$}
		
		\begin{scope}[-Latex]
		\path [-Latex] (0) edge node[above] {$12/4$} (1);
		\path [-Latex] (0) edge[bend right] node[above] {$12/4$} (2);
		\path [-Latex] (0) edge[bend left] node[above] {$8/4$} (4);
		\path [-Latex] (1) edge node[right] {$1/2$} (2);
		\path [-Latex] (1) edge node[above] {$11/2$} (3);
		\path [-Latex] (2) edge node[left] {$11/2$} (3);
		\path [Square-Latex] (3) edge node[above] {$2/1$} (4);
		\path [-Latex] (3) edge node[above] {$8/2$} (5);
		\path [-Latex] (4) edge node[above] {$8/1$} (6);
		\path [-Latex] (5) edge node[above] {$8/1$} (6);
		\end{scope}
		\end{tikzpicture}
		\subcaption{Ein schmaler Fluss mit Zurücksetzen auf $(d,e)$.}
		\label{subfig-thin-flow-with-resetting}
	\end{subfigure}
\caption{Ein schmaler Fluss einmal ohne Zurücksetzen und einmal mit Zurücksetzen auf $(d,e)$.
	In~\ref{subfig-thin-flow-without-resetting} und~\ref{subfig-thin-flow-with-resetting} hat die Kantenbeschriftung die Form $f_e/u_e$ und in den Knoten steht die Knotenauslastung.}
\end{figure}

\begin{lemma}\label{lemma-thin-flow-t-def}
	Ein $s$-Fluss $f$ in einem Netzwerk mit Zurücksetzen auf $E_1$ und Ursprung $s$ ist genau dann ein schmaler Fluss mit Zurücksetzen auf $E_1$, wenn eine Knotenbewertung $l\in\R_{\geq 0}^V$ existiert, die die folgenden Bedingungen erfüllt:
	\begin{enumerate}[label=(T\arabic*)]
		\item\label{def-thin-flow-source} $l_s = 0$,
		\item\label{def-thin-flow-x-zero} $l_w \leq l_v$, \tabto{4.5cm} für $vw\in E \setminus E_1$ mit $f_{vw}=0$,
		\item\label{def-thin-flow-x-positive} $l_w = \max\{ l_v, f_{vw} / u_{vw} \}$,  \tabto{4.5cm} für $vw\in E\setminus E_1$ mit $f_{vw} > 0$,
		\item\label{def-thin-flow-resetting-edge} $l_w = f_{vw} / u_{vw}$,  \tabto{4.5cm} für $vw\in E_1$,
		\item\label{def-thin-flow-no-resetting-edge} $l_w \geq \min_{vw\in \delta^-(w)} l_v$, \tabto{4.5cm} falls $\delta^-(w)\cap E_1 = \emptyset$.
	\end{enumerate}
	Diese Knotenbewertung stimmt dann mit der Knotenauslastung von $f$ überein.
	Außerdem gelten die Bedingungen~\ref{def-thin-flow-source}, \ref{def-thin-flow-x-zero} und~\ref{def-thin-flow-no-resetting-edge} bereits für die Knotenauslastungen von $f$, falls $f$ nur ein $s$-Fluss ist.
\end{lemma}
\begin{proof}
	Sei $f$ ein schmaler Fluss mit Zurücksetzen auf $E_1$ und sei $l\in\R_{\geq 0}^V$ die Knotenauslastung von $f$.
	Man zeige, dass $l$ die Bedingungen~\ref{def-thin-flow-source}-\ref{def-thin-flow-no-resetting-edge} erfüllt, und benutze dabei die Darstellung aus Proposition~\ref{prop-congestion-labels-dijkstra}.
	Wegen $l_s = 0$ gilt bereits~\ref{def-thin-flow-source}.
	Für eine Kante $vw\in E\setminus E_1$ mit $f_{vw}=0$ gilt $l_w\leq \rho_{vw}(l_v, f_{vw}) = l_v$, wodurch auch~\ref{def-thin-flow-x-zero} folgt.
	Ist $vw\in E$ mit $f_{vw} > 0$, so existiert ein $s$-$w$-Pfad mit positivem Fluss, der die Kante $vw$ benutzt, sodass die Auslastung dieses Pfades $l_w=\rho_{vw}(l_v, f_{vw})$ ist, da der $s$-$v$-Teilpfad die Auslastung $l_v$ besitzt.
	Für $vw\notin E_1$ bedeutet das gerade $l_w = \max\{ l_v, f_{vw}/u_{vw} \}$ also~\ref{def-thin-flow-x-positive}, und für $vw\in E_1$ zeigt das $l_w = f_{vw} / u_{vw}$.
	Für~\ref{def-thin-flow-resetting-edge} bleibt der Fall $f_{vw} = 0$ zu prüfen:
	Dann ist aber $l_w = \min_{uw\in\delta^-(w)} \rho_{uw}(l_u, f_{uw}) = f_{vw} / u_{vw} = 0$.
	Zuletzt betrachte man den Fall, dass $w$ keine eingehende zurücksetzende Kante hat.
	Dann folgt aber bereits $l_w = \min_{vw\in \delta^-(w)} \max\{ l_v, f_{vw} / u_{vw} \} \geq \min_{vw\in\delta^-(w)} l_v$, was auch~\ref{def-thin-flow-no-resetting-edge} impliziert.
	
	Sei nun umgekehrt $l\in\R_{\geq 0}^V$ eine Knotenbewertung, die~\ref{def-thin-flow-source}-\ref{def-thin-flow-no-resetting-edge} erfüllt.
	Man verwende eine Induktion über die Distanz eines Knotens $w$ zu $s$ bezüglich der Kantenzahl, um zu zeigen, dass $l_w$ die Knotenauslastung von  $f$ ist.
	Für $w=s$ gilt $l_s=0$ bereits nach~\ref{def-thin-flow-source}.
	Für $w\neq s$ ist $l_w$ nach~\ref{def-thin-flow-x-zero},~\ref{def-thin-flow-x-positive} und~\ref{def-thin-flow-resetting-edge} eine untere Schranke an $\rho_{vw}(l_v, f_{vw})$ für $vw\in\delta^-(w)$.
	Es bleibt zu zeigen, dass für eine eingehende Kante $vw$ der Wert $\rho_{vw}(l_v, f_{vw})$ auch $l_w$ annimmt:
	Falls eine eingehende Kante $vw\in E_1$ oder eine eingehende Kante $vw$ mit $f_{vw} > 0$ existiert, ist $\rho_{vw}(l_v, f_{vw}) = l_w$ nach~\ref{def-thin-flow-resetting-edge} und~\ref{def-thin-flow-x-positive}.
	Sonst ist $l_w\geq \min_{vw\in \delta^-(w)} l_v = \min_{vw\in \delta^-(w)} \rho_{vw}(l_v', x_{vw}')$ nach~\ref{def-thin-flow-no-resetting-edge}, was die Behauptung zeigt.
	Um zu sehen, dass $f$ nun ein schmaler Fluss mit Zurücksetzen auf $E_1$ ist, genügen die Bedingungen~\ref{def-thin-flow-source}, \ref{def-thin-flow-x-positive} und~\ref{def-thin-flow-resetting-edge}, da dadurch jeder $s$-$v$-Teilpfad eines Pfades mit positivem Fluss gerade Auslastung $l_v$ hat.
\end{proof}

\begin{remark}\label{remark-thin-flow}
	Hier ist, wie in~\cite[Definition~4]{Cominetti2011}, im Vergleich zu~\cite[Definition~6]{Koch2011} die Bedingung~\ref{def-thin-flow-no-resetting-edge} zusätzlich eingeführt worden.
	Diese ist nötig, um im Beweis von Theorem~\ref{thm-alpha-extension-is-nash-flow} zu zeigen, dass die erweiterten Ankunftszeiten tatsächlich mit den angegebenen übereinstimmen.
	Ohne diese Bedingung würde dies nämlich nicht gelten, wie man in Abbildung~\ref{figure-labels} erkennen kann.
\end{remark}
\todo{
	Introduce network inflow with normalized thin-flows
	\begin{definition}[Schmaler Fluss]\label{def-thin-flow}
		Der Fluss $x'$ heißt \emph{normiert}, falls $F=d$ gilt.
		Ist $E_1$ nichtleer, so heißt $x'$ \emph{schmaler Fluss mit Zurücksetzen auf $E_1$}, sonst heißt $x'$ \emph{schmaler Fluss ohne Zurücksetzen}.
	\end{definition}
}

\begin{lemma}\label{lemma-equivalent-thin-flow}
	Ein $s$-Fluss $f$ ist genau dann ein schmaler Fluss mit Zurück\-setzen auf $E_1$, wenn $f_{vw}= 0$ für alle $vw\in E$ mit $l_w < \rho_{vw}(l_v, f_{vw})$ gilt.
	Dabei sei $l\in\R_{\geq 0}^V$ die Knotenauslastung von $f$.
\end{lemma}
\begin{proof}
	Angenommen, $f$ sei ein schmaler Fluss mit Zurücksetzen auf $E_1$.
	Für eine Kante $vw\in E$ mit $l_w < \rho_{vw}(l_v, f_{vw})$ ist $vw\notin E_1$ und $f_{vw}=0$ nach~\ref{def-thin-flow-resetting-edge} und~\ref{def-thin-flow-x-positive}.

	Nun nehme man an, es gelte $f_{vw}=0$ für alle $vw\in E$ mit $l_w < \rho_{vw}(l_v, f_{vw})$, und man zeige die Eigenschaften~\ref{def-thin-flow-x-positive} und \ref{def-thin-flow-resetting-edge}.
	Der Rest folgt dann mit Lemma~\ref{lemma-thin-flow-t-def}.
	Ist $vw\notin E_1$ mit $f_{vw}>0$, so gilt $l_w = \rho_{vw}(l_v, f_{vw})$, wodurch Bedingung~\ref{def-thin-flow-x-positive} folgt.
	Für~\ref{def-thin-flow-resetting-edge} betrachte man zunächst $vw\in E_1$ mit $f_{vw}>0$: Hier folgt wieder $l_w = \rho_{vw}(l_v, f_{vw})$.
	Falls $f_{vw}=0$ gilt, ist $l_w=\min_{\tilde{v}w\in \delta^-(w)} \rho_{\tilde{v}w}(l_{\tilde{v}}, f_{\tilde{v}w}) = 0$.
\end{proof}

\todo{Vlt im folgenden Korollar auch/oder Gleichheit zu $\max f_e/u_e$ zeigen}
\begin{corollary}\label{cor-thin-flows-max-node-label-equals-max-edge-label}
	Für einen schmalen $s$-Fluss $f$ in einem Netzwerk mit nichtleerer Kantenmenge und zugehöriger Knotenauslastung $l$ gilt
	\[
		\max_{vw\in E} \rho_{vw}(l_v, f_{vw}) = \max_{v\in V} l_v
	\]
\end{corollary}
\begin{proof}
	Für $w\neq s$ ist $l_w = \min_{vw\in\delta^-(w)} \rho_{vw}(l_v, x_{vw})$.
	Daher gilt bereits \glqq$\geq$\grqq.
	Angenommen es gebe eine Kante $vw\in E$ mit $\rho_{vw}(l_v, f_{vw}) > \max_{u\in V} l_u$
	Dann ist aber bereits $f_{vw}/u_{vw} > \max_{u\in V} l_u \geq 0$, womit $l_w \geq \rho_{vw}(l_v, f_{vw})$ nach Lemma~\ref{lemma-equivalent-thin-flow} gelten müsste.
\end{proof}

\begin{lemma}[Eindeutigkeit der Knotenauslastung]\label{lemma-node-congestion-unique}
	Seien ein azyklisches Netzwerk mit Ursprung $s$ und Versorgungsrate $d\in\R_{\geq 0}$ sowie $E_1\subseteq E$ gegeben.
	Dann sind die Knotenauslastungen aller normierter schmaler Flüsse mit Zurücksetzen auf $E_1$ identisch.
\end{lemma}
\begin{proof}
	Angenommen, es existieren zwei normierte schmale Flüsse $x'$ und $y'$ mit Zurück\-setzen auf $E_1$ mit unterschiedlichen Knotenauslastungen $l'\neq h'$.
	Dann ist oBdA. die Menge $S:=\{ v\in V \mid l_v' > h_v' \}$ nichtleer.
	Da $x'$ und $y'$ beide $b$-Flüsse sind, gilt 
	\begin{equation}\label{proof-uniqueness-labels-eq}
	\sum_{e\in \delta^+(S)} x'_e - \sum_{e\in\delta^-(S)} x'_e = \sum_{v\in S} b_v = \sum_{e\in\delta^+(S)} y_e' - \sum_{e\in\delta^-(S)} y_e'.
	\end{equation}
	Für Kanten $vw\in \delta^+(S)$ gilt $x_e' \leq y_e'$, da für $x_e' > y_e'$ Lemma~\ref{lemma-equivalent-thin-flow} und $v\in S$ die Ungleichung $l_w' = \rho_{vw}(l_v', x_{vw}') > \rho_{vw}(h_v', y_{vw}')\geq h_w'$ implizieren würden, welche jedoch im Widerspruch zu $w\notin S$ steht.
	Ebenso gilt für Kanten $vw\in\delta^-(S)$ die Ungleichung $x_e' \geq y_e'$, weil sonst $l_w' \leq \rho_{vw}(l_v', x_{vw}') \leq \rho_{vw}(h_v', y_{vw}') = h_w'$ wegen $v\notin S$ gelten würde, was $w\in S$ widerspricht.
	Gleichung~\ref{proof-uniqueness-labels-eq} impliziert dann sogar $x_e = y_e$ für alle $e\in \delta(S):=\delta^-(S) \cup \delta^+(S)$.
	Für $vw\in \delta^-(S)$ ist dann $y_e'=x_e'=0$, da für $y_e'=x_e'>0$ wieder mit $v\notin S$ die Ungleichung $l_w' = \rho_{vw}(l_v, x_{vw}')\leq \rho_{vw}(h_v, y_{vw}')=h_w'$ der Bedingung $w\in S$ widerspricht.
	
	Aufgrund der Azyklizität existiert ein Knoten $w\in S$ mit $\delta^-(w)\subseteq \delta^-(S)$.
	Eingehende Kanten von $w$ sind daher nicht in $E_1$ enthalten, da für solch eine Kante $\rho_{vw}(l_v', x_e') = 0 = \rho_{vw}(l_v', y_e')$ gelten würde, sodass $l_w' = 0 = h_w'$ der Voraussetzung $w\in S$ widerspricht.
	Daher ist $l_w' = \min_{vw\in \delta^-(w)} l_v'$ und $h_w' = \min_{vw\in\delta^-(w)} h_v'$, was den Widerspruch $l_w' \leq h_w'$ impliziert.
\end{proof}

Zwar sind schmale Flüsse im Allgemeinen nicht eindeutig -- nichtmal ohne zurück\-setzende Kanten --, wie man leicht in Abbildung~\ref{figure-thin-flows-not-unique} erkennen kann, jedoch ist der Fluss auf einigen Kanten eindeutig:

\begin{figure}
	\centering
	\begin{subfigure}{.45\textwidth}
		\begin{tikzpicture}[scale=0.8]
		\node[draw, circle] (S) at (0,2) {$s$};
		\node[draw, circle] (U) at (2,2) {$u$};
		\node[draw, circle] (W) at (4,1) {$w$};
		\node[draw, circle] (V) at (4,3) {$v$};
		\node[draw, circle] (T) at (6,2) {$t$};
		
		\path [->] (S) edge node[above] {$2$} (U);
		\path [->] (U) edge node[above] {$1$} (V);
		\path [->] (U) edge node[below] {$1$} (W);
		\path [->] (V) edge node[above] {$1$} (T);
		\path [->] (W) edge node[below] {$1$} (T);
		\end{tikzpicture}
	\end{subfigure}
	\begin{subfigure}{.45\textwidth}
		\begin{tikzpicture}[scale=0.8]
		\node[draw, circle] (S) at (0,2) {$s$};
		\node[draw, circle] (U) at (2,2) {$u$};
		\node[draw, circle] (W) at (4,1) {$w$};
		\node[draw, circle] (V) at (4,3) {$v$};
		\node[draw, circle] (T) at (6,2) {$t$};
		
		\path [->] (S) edge node[above] {$2$} (U);
		\path [->] (U) edge node[above] {$2$} (V);
		\path [->] (U) edge node[below] {$0$} (W);
		\path [->] (V) edge node[above] {$2$} (T);
		\path [->] (W) edge node[below] {$0$} (T);
		\end{tikzpicture}
	\end{subfigure}
	\caption{Zwei verschiedene schmale Flüsse derselben Instanz ohne Zurück\-setzen. Auf den Kanten steht der Flusswert; die Kapazitäten sind allesamt $1$.}
	\label{figure-thin-flows-not-unique}
\end{figure}

\begin{corollary}[Eindeutigkeit schmaler Flüsse]
	 Für eine Instanz von~$\problemThinFlow$ sind alle Lösungen auf den zurück\-setzen\-den Kanten sowie auf Kanten $vw$ mit $l_v \neq l_w$ identisch.
	 Dabei ist $(l_v)_{v\in V}$ die eindeutige Knotenauslastung.
\end{corollary}
\begin{proof}
	Da nach Lemma~\ref{lemma-node-congestion-unique} die Knotenauslastung $(l_v)_{v\in V}$ für jeden schmalen Fluss $f$ eindeutig ist, gilt $f_{vw} = u_{vw} l_w$ für $vw\in E_1$ nach~\ref{def-thin-flow-resetting-edge}.
	Für $vw\notin E_1$ betrachte man zunächst den Fall $l_v < l_w$.
	Nach Definition gilt $l_w \leq \max \{ l_v, f_{vw}/u_{vw} \}$, sodass man $f_{vw}/u_{vw} > l_v \geq 0$ folgern kann.
	Nach Lemma~\ref{lemma-equivalent-thin-flow} gilt also $f_{vw} = l_w u_{vw}$.
	Im Fall $l_v > l_w$ muss $f_{vw} = 0$ gelten, da sonst $l_w=\max\{ l_v, f_e/u_e \}$ gelten würde.
\end{proof}

\begin{theorem}[Existenz schmaler Flüsse]\label{thm-existence-thin-flow}
	Für jede Instanz von $\problemThinFlow$ existiert ein schmaler Fluss, falls im zugehörigen Netzwerk ein $s$-Fluss existiert.
\end{theorem}

Um die Existenz zu beweisen, benötigen wir zunächst den Fixpunktsatz von Kakutani.
Dazu führen wir die folgenden Begriffe ein:

\begin{definition}[Korrespondenz, Fixpunkt]
	Eine \emph{Korrespondenz} von einer Menge $A$ in eine Menge $B$ ist eine Abbildung $f: A \to \mathcal{P}(B)\setminus \{ \emptyset \}$ von $A$ in die Potenzmenge von $B$ ohne die leere Menge.
	Sind $A$ und $B$ topologische Räume, so nennt man $f$ \emph{abgeschlossen}, wenn die zugehörige Relation $R_f := \{ (a,b) \in A\times B \mid b\in f(a)  \}$ in der Produkttopologie abgeschlossen ist.

	Ein \emph{Fixpunkt} einer Korrespondenz $f: X \to \mathcal{P}$ ist ein Punkt $x\in X$ mit $x\in f(x)$.
\end{definition}

Nun lautet der Fixpunktsatz von Kakutani (siehe~\cite{Heuser1991Fix}):

\begin{satz}[Fixpunktsatz von Kakutani]\label{satz-kakutani}
	Seien $C\subseteq E$ eine nichtleere, konvexe und kompakte Teilmenge eines normierten Raumes $E$ und $f: C \to \mathcal{P}(C)$ eine abgeschlossene, konvexwertige Korrespondenz.
	Dann besitzt $f$ einen Fixpunkt.
\end{satz}

\begin{proof}[Beweis von Theorem~\ref{thm-existence-thin-flow}]
	Man betrachte die Menge $C$ aller $b$-Flüsse als Teilmenge des metrischen Raums $\R^E$.
	Nach Voraussetzung ist diese nichtleer, konvex, da für $f, g\in C, \lambda \in [0,1]$ der Fluss $\lambda f + (1-\lambda)g$ wieder die Flusserhaltung erfüllt, und kompakt, da jeder Fluss in $C$ aufgrund der Azyklizität in $s$-$v_i$-Wege zerlegbar ist, wodurch jede Kante maximal Fluss $b_s$ besitzen kann, sodass $C$ eine beschränkte Teilmenge von $\subseteq [0, b_s]^E$ ist.
	Außerdem ist $C$ abgeschlossen, da es die Lösungsmenge eines linearen Gleichungssystem ist.
	Auf $C$ definiere man die Korrespondenz
	\[
	\Gamma: C\to \mathcal{P}(C)\setminus\{ \emptyset \}, ~ f\mapsto \{ g\in C \mid \forall vw\in E: l_w < \rho_{vw}(l_v, f_{vw}) \implies g_{vw} = 0 \}.
	\]
	Dabei sei $l\in\R_{\geq 0}^V$ die zu $f$ gehörige Knotenauslastung.
	Diese Korrespondenz ist wohldefiniert, da $\Gamma(f)$ für jedes $f\in C$ nichtleer ist:
	Der Knoten $s$ kann im Graphen $G':= (V, E')$ mit $E':= \{ vw\in E \mid \rho_{vw}(l_v, f_{vw})=l_w \}$ jeden Knoten erreichen, da $G'$ azyklisch ist und jeder Knoten $w\neq s$ mindestens eine eingehende Kante mit $l_w = \rho_{vw}(l_v, f_{vw})$ besitzt. 
	Daher existiert ein $s$-$v$-Weg $P_v$ in $G'$ für alle $v\in V$ und der Fluss $\sum_{v\in V\setminus\{ s\}} b_v \cdot P_v$ ist in $\Gamma(f)$ enthalten.
	
	Für jedes $f\in C$ ist $\Gamma(f)$ konvex: Ist für $g, h\in \Gamma(f), \lambda\in [0,1]$ und eine Kante $vw\in E$ die Bedingung $l_w < \rho_{vw}(l_w, f_{vw})$ erfüllt, so gilt auch $\lambda g_{vw} + (1-\lambda) h_{vw} = 0$, womit $\lambda g + (1-\lambda)h$ im Bild $\Gamma(f)$ liegt.
	Des Weiteren ist $\Gamma$ abgeschlossen: 
	Sei eine konvergente Folge $((f^{n}, g^{n}))_{n\in\N}$ in $R_\Gamma$ mit Grenzwert $(f, g)$ gegeben.
	Die zu $f$ bzw. $f^n$ gehörigen Knotenauslastungen seien gegeben durch $l$ bzw. $l^n$.
	Angenommen, es gelte $l_w < \rho_{vw}(l_v, f_{vw})$.
	Da die Zuordnung $x'\mapsto l'$ eines $s$-Flusses auf seine Knotenauslastung nach Lemma~\ref{prop-congestion-labels-dijkstra} stetig ist, existiert ein $N\in\N$ mit $l_w^n<\rho_{vw}(l_v^n, f_{vw}^n)$ für alle $n\geq N$.
	Damit ist auch $g_{vw} = \lim_{n\to\infty} g_{vw}^n = 0$.
	
	Daher existiert nach dem Fixpunktsatz von Kakutani~\ref{satz-kakutani} ein Fixpunkt von $\Gamma$, welcher nach Lemma~\ref{lemma-equivalent-thin-flow} ein schmaler Fluss mit Zurücksetzen auf $E_1$ ist.
\end{proof}

\section{Effiziente Berechnung schmaler Flüsse}

Die Existenz schmaler Flüsse mit Zurücksetzen aus Theorem~\ref{thm-existence-thin-flow} suggeriert bereits einen exponentiellen Algorithmus, um einen solchen Fluss mit Wert $d$ zu berechnen:
Man rate die Menge $E'\subseteq E$ aller Kanten $vw$, die $l_w' = \rho_{vw}(l_v', x_{vw}')$ erfüllen, sowie für alle $vw\in E'\setminus E_1$ eine Binärzahl $\sigma_{vw}\in\{0,1\}$, die $\max\{ l_v', x_{vw}/u_{vw}\} = \sigma_{vw} l_v' + (1-\sigma_{vw}) x_{vw}/u_{vw}$ erfüllt.
Dann suche man einen zulässigen Punkt in folgendem Polyeder:
\begin{align*}
&	(x', l') \in \R_{\geq 0}^{E'\times V} \\
	\text{u.d.N.}\quad &
	x'(\delta^+(v)) - x'(\delta^-(v)) =
		\begin{cases}
			d, & \text{falls $v=s$},\\
			0, & \text{sonst,} 
		\end{cases} & \text{für $v\in V\setminus\{t \}$,} \\
	& l'_s = 0, \\
	& \begin{array}{@{}l}
		l_w' = \sigma_{vw} l_v' + (1-\sigma_{vw}) x'_{vw}/u_{vw}\\
		\sigma_{vw} l_v' + (1-\sigma_{vw}) x'_{vw}/u_{vw} \geq x'_{vw}/u_{vw}\\
		\sigma_{vw} l_v' + (1-\sigma_{vw}) x'_{vw}/u_{vw} \geq l_v'
	\end{array}
	& \text{für $vw\in E'\setminus E_1$},\\
	& l_w' = x'_{vw}/u_{vw} & \text{für $vw\in E_1$.}
\end{align*}
\todo{Genauerer Beweis}

\newcommand{\TFNP}{\mathbf{TFNP}}

Existiert solch ein Punkt $(x', l')$, so ist $x'$ nach Lemma~\ref{lemma-equivalent-thin-flow} ein schmaler Fluss mit Zurücksetzen auf $E_1$ mit zugehöriger Knotenbewertung $l'$.
Außerdem muss aufgrund der Existenz eines schmalen Flusses mindestens eine Wahl von $E'$ und $\sigma$ existieren, sodass der zugehörige Polyeder nichtleer ist.
\todo{Habe ich hier nicht das Problem, dass ich das nur für $\Q$ in polynomieller Zeit berechnen kann?}
Die Zulässigkeit eines solchen Polyeders kann in polynomieller Zeit beispielsweise durch die Ellipsoidmethode ermittelt werden. 
Daher liegt das Problem, zu einem Netzwerk $(V, E, u, s, t)$ einen schmalen $d$-Fluss mit Zurücksetzen auf $E_1\subseteq E$ zu finden, in der Komplexitätsklasse $\TFNP$:

\begin{definition}[Komplexitätsklasse $\TFNP$]
	Eine Relation $P(x,y)$ liegt in der Komplexitätsklasse $\TFNP$ genau dann, wenn es einen deterministischen, polynomiellen Algorithmus gibt, der entscheidet, ob $P(x,y)$ für zwei Kandidaten $x$ und $y$ gilt, und wenn für alle $x$ ein Kandidat $y$ existiert, dessen Kodierung polynomielle Größe in der Kodierungslänge von $x$ besitzt, sodass $P(x, y)$ gilt.
\end{definition}

In diesem Abschnitt wird ein polynomieller Algorithmus zur Berechnung von schmalen Flüssen ohne Zurücksetzen vorgestellt.

\begin{lemma}\label{lemma-thin-flows-without-resetting-have-minimal-congestion}
	Sei ein schmaler $b$-Fluss $f$ ohne Zurücksetzen mit zugehöriger Knotenauslastung $l$ gegeben.
	Dann ist $f$ ein $b$-Fluss minimaler Auslastung $q$ und für $E\neq\emptyset$ ist $X:=\{ v\in V \mid l_v < q \}$ ein dünnster Schnitt.
\end{lemma}
\begin{proof}
	Ist $b$ der Nullvektor -- wie im Fall $E=\emptyset$ --, ist $f$ der Nullfluss, da das Netzwerk azyklisch ist.
	
	Sei also $q>0$ die Auslastung von $f$ im Falle von $b\neq 0$.
	Nach Korollar~\ref{cor-thin-flows-max-node-label-equals-max-edge-label} genügt es zu
	zeigen, dass die Auslastung jeder ausgehenden Kante von $X$ bezüglich $f$ gerade $q$ beträgt und auf eingehenden Kanten von $X$ kein Fluss fließt.
	Wegen $l_s=0$ ist $s$ im Schnitt $X$ enthalten.
	Sei $vw\in\delta^+(X)$ eine ausgehende Kante.
	Dann gilt $l_v < l_w = q $ und nach Proposition~\ref{prop-congestion-labels-dijkstra} auch $q = l_w \leq \rho_{vw}(l_v, f_{vw}) = f_{vw} / u_{vw}$.
	Da $q$ eine obere Schranke an $f_{vw}/u_{vw}$ ist, gilt also $f_{vw}/u_{vw} = q$.
	Für eine eingehende Kante $vw\in\delta^-(X)$ gilt $\rho_{vw}(l_v, f_{vw})\geq l_v > l_w$ und somit muss nach Lemma~\ref{lemma-equivalent-thin-flow} der Fluss $f$ auf der Kante $vw$ verschwinden.
\end{proof}

\begin{definition}
	Seien $\mathcal{I}:=(V, E, u, b)$ eine $b$-Fluss-Instanz und $X\subseteq V$ ein dünnster Schnitt mit Auslastung $q$.
	Dann heißt $\mathcal{I}[X] := (X, E[X], u[X], b[X])$ die \emph{durch $X$ induzierte Teilinstanz von $\mathcal{I}$},
	wobei $E[X] := \{ vw \in E \mid v, w \in X \}$ als die Menge der Kanten zwischen Knoten in $X$, $u[X]$ als die Einschränkung von $u$ auf $E[X]$ und $b[X]$ durch \[
	b[X]_v := b_v - u(\delta^+_E(v)\cap\delta^+_E(X)) q \quad \text{für $v\in X$}
\] definiert ist.
\end{definition}

Man bemerke, dass $b[X]_v \leq b_v$ für $v\in X$ erfüllt und $b[X]$ ein zulässiger Balancevektor ist: \[
	\sum_{v\in X} b[X]_v = b(X) - u(\delta^+_E(X)) q = 0
\]

\begin{proposition}\label{prop-restricted-minimal-flow-is-b-flow-on-induced-instance}
	Ist $f$ ein $b$-Fluss minimaler Auslastung einer $b$-Fluss-Instanz $\mathcal{I}$, so ist die Einschränkung $f':=f_{\mid E[X]}$ ein $b[X]$-Fluss der durch $X$ induzierten Instanz $(X, E[X], u[X], b[X])$ für jeden dünnsten Schnitt $X$.
	Ist $f$ sogar ein $s$-Fluss mit Knotenauslastung $l$ und $b_s > 0$, so stimmt die Einschränkung von $l$ auf $X$ mit der Knotenauslastung $l'$ vom $s$-Fluss $f'$ überein.
\end{proposition}
\begin{proof}
	Nach Korollar~\ref{cor-easy-characterization-sparsest-cut} gelten $f_{e}/u_{e}=q$ für $e\in\delta^+(X)$ und $f_{e} = 0$ für $e\in\delta^-(X)$.
	Daher kann man zeigen, dass $f'$ ein $b[X]$-Fluss ist, denn für alle $v\in X$ gilt:
	\begin{align*}
	f'(\delta^+_{E[X]}(v)) - f'(\delta^-_{E[X]}) &= f(\delta^+_E(v)) - f(\delta^-_E(v)) - f(\delta^+_{E[X]}(v)\cap\delta^+_{E[X]}(X)) \\
	&= b_v - u(\delta^+_{E[X]}(v)\cap\delta^+_{E[X]}(X))q = b[X]_v.
	\end{align*}
	Sei nun $f$ sogar ein $s$-Fluss mit $b_s > 0$, wodurch $s$ in $X$ enthalten ist.
	Sei $v\in X$ gegeben.
	Da jeder $s$-$v$-Pfad in $(X, E')$ auch in $(V, E)$ enthalten ist, gilt bereits $l'_v\leq l_v$.
	Außerdem liegt jeder $s$-$v$-Pfad in $(V, E)$ mit positivem Fluss bezüglich $f$ auch komplett in $(X, E')$, da $f$ auf $\delta^-_E(X)$ verschwindet.
	Daraus folgt $l'_v \geq l_v$.
\end{proof}

Daher kann man bereits erkennen, dass ein schmaler Fluss ohne Zurücksetzen eingeschränkt auf eine durch einen dünnsten Schnitt induzierte Teilinstanz wieder ein schmaler Fluss ist.
Das Ziel ist es nun, aus einem schmalen Fluss der Teilinstanz einen schmalen Fluss auf der ursprünglichen Instanz zu konstruieren.

\begin{lemma}\label{lemma-flow-minimal-congestion-sparsest-cut-then-outside-X-congestion-q}
	Seien ein $s$-Fluss $f$ minimaler Auslastung $q^*$ mit zugehöriger Knotenauslastung $l$ und ein dünnster Schnitt $X$ in einem Netzwerk ohne zurücksetzende Kanten gegeben.
	Für alle $v\in V\setminus X$ hat jeder $s$-$v$-Pfad $P$ Auslastung $q^*=l_v$.
\end{lemma}
\begin{proof}
	Ist $b$ der Nullvektor, so ist $f$ der Nullfluss und $l_v = 0$ für alle $v\in V$.
	Sonst ist $s$ in $X$ enthalten und für alle $v\in V\setminus X$ enthält jeder $s$-$v$-Pfad mindestens eine Kante, die $X$ verlässt und nach Korollar~\ref{cor-easy-characterization-sparsest-cut} Auslastung $q^*$ besitzt.
	Da das Netz keine zurücksetzenden Kanten hat, beträgt die Auslastung also $\max_{e\in P} f_e / u_e = q^*$.
\end{proof}

\begin{corollary}\label{cor-thin-flow-alg-correct}
	Seien ein $s$-Fluss $f$ minimaler Auslastung sowie ein dünnster Schnitt $X$ auf einem Netzwerk ohne zurücksetzende Kanten gegeben.
	Ist $f'$ ein schmaler Fluss ohne Zurücksetzen auf der durch $X$ induzierten Teilinstanz, so ist $g := f' \sqcup f$ definiert durch \[
		g_e := \begin{cases}
			f'_e, & \text{falls $e\in E[X]$,} \\
			f_e, & \text{falls $e\in E\setminus E[X]$,}
		\end{cases}
	\] ein schmaler Fluss ohne Zurücksetzen auf der ursprünglichen Instanz.
\end{corollary}
\begin{proof}
	Es ist leicht zu sehen, dass $g$ ein $s$-Fluss minimaler Auslastung ist:
	Für einen Knoten $v \in V\setminus X$ ist $g(\delta^+_E(v)) - g(\delta^-_E(v)) = f(\delta^+_E(v)) - g(\delta^-_E(v)) = b_v$ und für $v\in X$ gilt
	\begin{align*}
	g(\delta^+_E(v)) - g(\delta^-_E(v)) =& f'(\delta^+_{E[X]}(v))  - f'(\delta^-_{E[X]}(v))  \\
	&+ f(\delta^+_E(v)\cap\delta^+_E(X)) - f(\delta^-_E(v)\cap\delta^-_E(X)) = b_v
	\end{align*}

	Sei ein $s$-$v$-Pfad $P$ mit positivem Fluss gegeben.
	Ist $v$ nicht in $X$, so hat $P$ nach Lemma~\ref{lemma-flow-minimal-congestion-sparsest-cut-then-outside-X-congestion-q} Auslastung $l_v$.
	Ist $v$ in $X$ enthalten, muss $P$ entlang Kanten innerhalb von $X$ laufen, da $g$ auf eingehenden Kanten von $X$ verschwindet.
	Nach Proposition~\ref{prop-restricted-minimal-flow-is-b-flow-on-induced-instance} stimmt die Auslastung von $v$ bezüglich der Einschränkung von $f$ mit der Auslastung von $v$ bezüglich $f$ überein, sodass auch die Auslastung von $P$ mit der Auslastung von $v$ bezüglich $f$ übereinstimmen.
\end{proof}

\begin{algorithm}
\caption{Berechnung eines schmalen Flusses ohne Zurücksetzen}
\label{algorithm-computation-thin-flow-without-resetting}
\begin{algorithmic}[1]
	\Procedure{ThinFlow}{$V, E, u, b, s$}
	\State $f \gets 0^E$
	\State $(V', E', b') \gets (V, E, b)$
	\While{$b_s \neq 0$}
	\State Berechne auslastungsminimalen $s$-Fluss $f'$ und dünnsten Schnitt $X$
	\Statex auf $(V', E', b')$ mit Kapazitäten $u_{\mid E'}$; melde Fehler, falls $f'$ nicht existiert.
	\State $f_{\mid E'\setminus E'[X]} \gets f'_{\mid E'\setminus E'[X]}$
	\State $(V', E', b') \gets (V'[X], E'[X], b'[X])$
	\EndWhile
	\State\Return $f$
	\EndProcedure
\end{algorithmic}
\end{algorithm}

Dies bildet die Grundlage des polynomiellen Algorithmus'~\ref{algorithm-computation-thin-flow-without-resetting} zur Berechnung von schmalen Flüssen ohne Zurücksetzen.

\begin{theorem}
		Algorithmus~\ref{algorithm-computation-thin-flow-without-resetting} berechnet einen schmalen Fluss ohne Zurücksetzen oder bricht mit einer Fehlermeldung ab, falls kein $b$-Fluss existiert.
		Er benötigt eine Laufzeit von $\bigO((\size{b} + \size{u})^2 n^3 m)$.
\end{theorem}
\begin{proof}
	Für die Korrektheit des Algorithmus benutzt man Korollar~\ref{cor-thin-flow-alg-correct}:
	Man prüfe die folgende Invariante der Schleife:
	\glqq $f$ verschwindet auf $E'$ und ist $g$ ein schmaler Fluss auf $(V',E',b')$, so ist $g \sqcup f$ ein schmaler Fluss auf $(V, E, b)$.\grqq\ 
	Zu Beginn gilt dies offenbar, da in diesem Fall $(V', E', b')$ mit $(V, E, b)$ übereinstimmt und $f$ mit Nullen initialisiert wurde.
	Nun betrachte man die Invariante im Verlauf einer Iteration.
	Seien dazu $(f^1, V^1, E^1, b^1)$ die Werte von $(f, V', E', b')$ zu Beginn der Iteration, für die die Invariante bereits gilt, und $(f^2, V^2, E^2, b^2)$ die Werte nach Ende der Iteration.
	Dann sei $f'$ der berechnete auslastungsminimale Fluss auf $(V^1, E^1, b^1)$ und $(V^2, E^2, b^2)$ die durch den dünnsten Schnitt $X$ induzierte Teilinstanz von $(V^1, E^1, b^1)$.
	Man nehme nun an, $g$ sei ein schmaler Fluss auf $(V^2, E^2, b^2)$.
	Dann ist $g\sqcup f^2_{\mid E^1} = g \sqcup f'$ nach Korollar~\ref{cor-thin-flow-alg-correct} ein schmaler Fluss auf $(V^1, E^1, b^1)$ und aufgrund der Invariante ist $g \sqcup f^2 = (g \sqcup f') \sqcup f^1$ ein schmaler Fluss auf $(V, E, b)$.
	Außerdem verschwindet $f^2$ auf $E^2$, da $f$ nur auf $E^1\setminus E^2$ verändert wird.
	Terminiert die Schleife, so ist $b_s=0$, wodurch die Voraussetzung der Invariante
	für den Nullfluss $f_{\mid E'}$ erfüllt ist, sodass $f$ ein schmaler Fluss ist.
	
	Existiert kein $b$-Fluss auf $(V, E)$, so existiert auch kein auslastungsminimaler Fluss und der Algorithmus in Zeile~$5$ bricht mit einer Fehlermeldung ab.
	Sonst terminiert der Algorithmus nach mindestens $n$ Iterationen, denn in jeder Iteration wird aus $V'$ mindestens ein Knoten entfernt.
	Außerdem kann ein auslastungsminimaler $s$-Fluss nach Theorem~\ref{thm-compute-minimal-con-flow} in $\bigO((\size{b}+\size{u})^2 n^2 m)$ Zeit berechnet werden.
	Daraus kann man nach~\ref{lemma-calc-sparsest-cut} in $\bigO(n+m)$ arithmetischen Operationen einen dünnsten Schnitt gewinnen, wodurch eine Gesamtlaufzeit von $\bigO((\size{b} + \size{u})^2 n^3 m)$ entsteht.
\end{proof}


\begin{example}
	
	\begin{figure}
		\newcommand{\newnode}[4]{\node[lul] (#1) at #2 {#3 \\[0.5em] #4};}
		\centering
		\begin{subfigure}{\textwidth}
			\begin{tikzpicture}[lul/.style={draw,
				ellipse,
				align=center,
				inner sep=0pt,
				outer sep=4pt,
				text width=7mm,
				minimum height=1.5cm
			},
			scale=1]
			
			\newnode{0}{(0,2)}{$a$}{$32$}
			\newnode{1}{(3,2)}{$b$}{$0$}
			\newnode{2}{(5,0)}{$c$}{$-2$}
			\newnode{3}{(7,2)}{$d$}{$-12$}
			\newnode{4}{(10,3)}{$e$}{$-2$}
			\newnode{5}{(10,1)}{$f$}{$0$}
			\newnode{6}{(13,2)}{$g$}{$-16$}
			
			\begin{scope}[-Latex]
			\path [-Latex] (0) edge node[above] {$10/4$} (1);
			\path [-Latex] (0) edge[bend right] node[above] {$4/4$} (2);
			\path [-Latex] (0) edge[bend left] node[above] {$18/4$} (4);
			\path [-Latex] (1) edge node[right] {$0/2$} (2);
			\path [-Latex] (1) edge[snake=snake,segment amplitude=.4mm,segment length=2mm,line after snake=2mm] node[above] {$16/2$} (3);
			\path [-Latex] (2) edge node[left] {$4/2$} (3);
			\path [-Latex] (3) edge node[above] {$0/1$} (4);
			\path [-Latex] (3) edge node[above] {$8/2$} (5);
			\path [-Latex] (4) edge[snake=snake,segment amplitude=.4mm,segment length=2mm,line after snake=2mm] node[above] {$8/1$} (6);
			\path [-Latex] (5) edge[snake=snake,segment amplitude=.4mm,segment length=2mm,line after snake=2mm] node[above] {$8/1$} (6);
			\end{scope}
			
			\begin{scope}
			\draw[dashed, color=blue] (11.5,-1) -- (11.5, 4.5);
			\node[color=blue] at (11,-0.5) {$X$};
			\end{scope}
			\end{tikzpicture}
			\subcaption{Erste Iteration. Die Auslastung der Flaschenhalskanten beträgt $8$.}
		\end{subfigure}
		\par\bigskip
		\begin{subfigure}{\textwidth}
			\begin{tikzpicture}[lul/.style={draw,
				ellipse,
				align=center,
				inner sep=0pt,
				outer sep=4pt,
				text width=7mm,
				minimum height=1.5cm
			},
			scale=1]
			
			\newnode{0}{(0,2)}{$a$}{$32$}
			\newnode{1}{(3,2)}{$b$}{$0$}
			\newnode{2}{(5,0)}{$c$}{$-2$}
			\newnode{3}{(7,2)}{$d$}{$-12$}
			\newnode{4}{(10,3)}{$e$}{$-10$}
			\newnode{5}{(10,1)}{$f$}{$-8$}
			
			\begin{scope}[-Latex]
			\path [-Latex] (0) edge node[above] {$10/4$} (1);
			\path [-Latex] (0) edge[bend right] node[above] {$12/4$} (2);
			\path [-Latex] (0) edge[bend left] node[above] {$10/4$} (4);
			\path [-Latex] (1) edge node[right] {$0/2$} (2);
			\path [-Latex] (1) edge[snake=snake,segment amplitude=.4mm,segment length=2mm,line after snake=2mm] node[above] {$10/2$} (3);
			\path [-Latex] (2) edge[snake=snake,segment amplitude=.4mm,segment length=2mm,line after snake=2mm] node[left] {$10/2$} (3);
			\path [-Latex] (3) edge node[above] {$0/1$} (4);
			\path [-Latex] (3) edge node[above] {$8/2$} (5);
			\end{scope}
			
			\begin{scope}
			\draw[dashed, color=blue, rounded corners=5mm] (6.75,-1) -- (5.5, 3.75) -- (12, 1.25);
			\node[color=blue] at (6.25,-0.5) {$X$};
			\end{scope}
			\end{tikzpicture}
			\subcaption{Zweite Iteration. Die Auslastung der Flaschenhalskanten beträgt $5$.}
		\end{subfigure}
		\par\bigskip
		\begin{subfigure}{.45\columnwidth}
			\begin{tikzpicture}[lul/.style={draw,
				ellipse,
				align=center,
				inner sep=0pt,
				outer sep=4pt,
				text width=7mm,
				minimum height=1.5cm
			}]
			
			\newnode{0}{(0,2)}{$a$}{$32$}
			\newnode{1}{(3,2)}{$b$}{$-10$}
			\newnode{2}{(5,0)}{$c$}{$-12$}
			\newnode{4}{(5,3)}{$e$}{$-10$}
			
			\begin{scope}[-Latex]
			\path [-Latex] (0) edge[snake=snake,segment amplitude=.4mm,segment length=2mm,line after snake=2mm] node[above] {$11/4$} (1);
			\path [-Latex] (0) edge[bend right,decorate, decoration=snake, segment amplitude=.4mm,segment length=2mm,line after snake=2mm] node[above] {$11/4$} (2);
			\path [-Latex] (0) edge[bend left] node[above] {$10/4$} (4);
			\path [-Latex] (1) edge node[right] {$1/2$} (2);

			\end{scope}
			
			\begin{scope}
			\draw[dashed, color=blue, rounded corners=5mm] (0, 0) -- (3, 3.5) -- (6, 0.5);
			\node[color=blue] at (0,0.6) {$X$};
			\end{scope}
			\end{tikzpicture}
			\subcaption{Dritte Iteration. Die Auslastung der Flaschenhalskanten beträgt $2{,}75$.}
		\end{subfigure} \hfill
		\begin{subfigure}{.45\columnwidth}
			\centering
			\begin{tikzpicture}[lul/.style={draw,
				ellipse,
				align=center,
				inner sep=0pt,
				outer sep=4pt,
				text width=7mm,
				minimum height=1.5cm
			},
			scale=1]
			
			\newnode{0}{(0,2)}{$a$}{$10$}
			\newnode{4}{(3,2)}{$e$}{$-10$}
			
			\begin{scope}[-Latex]
			\path [-Latex] (0) edge[snake=snake,segment amplitude=.4mm,segment length=2mm,line after snake=2mm] node[above] {$10/4$} (4);
			
			\end{scope}
			
			\begin{scope}
			\draw[dashed, color=blue, rounded corners=5mm] (1.5, 0.5) -- (1.5, 3.5);
			\node[color=blue] at (1,1) {$X$};
			\end{scope}
			\end{tikzpicture}
			\subcaption{Vierte Iteration. Die Auslastung der Flaschenhalskante beträgt $2{,}5$.}
		\end{subfigure}
		
		\caption{Algorithmus zur Bestimmung von schmalen Flüssen ohne Zurück\-setzen. Geschlängelte Kanten sind Flaschenhalskanten in den jeweiligen auslastungs\-minimalen Flüssen. Kanten $e$ sind mit $f_e/u_e$ und Knoten mit der Bezeichnung und der Balance beschriftet.}
	\end{figure}	
\end{example}

\chapter{Komplexitätsklasse PPAD}\label{sec-complexity-class-ppad}




\clearpage          % neue Seite für Literaturverzeichnis
\nocite*
\thispagestyle{empty}
\bibliography{literature}
\end{document}		
